\documentclass[12pt,a4paper]{article}
% ------------------------------------------------------------------ %
% Shared Persian (XePersian) preamble.  Compile with:  xelatex        %
% ------------------------------------------------------------------ %
\usepackage{amsmath,amssymb,amsthm}
\usepackage{graphicx}
\usepackage{booktabs}
\usepackage{array}
\usepackage{multirow}
\usepackage{longtable}
\usepackage[table,dvipsnames]{xcolor}
\usepackage{tcolorbox}
\usepackage{enumitem}
\usepackage[margin=2.4cm]{geometry}
\usepackage{caption}
\usepackage[hidelinks]{hyperref}

\tcbuselibrary{skins,breakable}

\newcommand{\userfonts}{C:/Users/victus/AppData/Local/Microsoft/Windows/Fonts/}

\usepackage{xepersian}                      % must be loaded LAST
\settextfont[Path=\userfonts, Extension=.ttf, Scale=1.0]{Yas}
\setlatintextfont[Scale=0.92]{Latin Modern Roman}

\theoremstyle{definition}
\newtheorem{theorem}{قضیه}[section]
\newtheorem{lemma}[theorem]{لم}
\newtheorem{proposition}[theorem]{گزاره}
\newtheorem{corollary}[theorem]{نتیجه}
\newtheorem{definition}[theorem]{تعریف}
\newtheorem{example}[theorem]{مثال}
\newtheorem{remark}[theorem]{نکته}
\renewcommand{\proofname}{اثبات}

\newcommand{\sig}{\sigma}
\newcommand{\ehat}{\hat{e}}
\newcommand{\mhat}{\hat{m}}
\newcommand{\rbar}{\bar{\rho}}
\newcommand{\mbar}{\bar{m}}
\newcommand{\code}[1]{\lr{\texttt{#1}}}

\newtcolorbox{answer}[1][]{
  breakable, colback=black!4, colframe=black!45, boxrule=0.5pt,
  left=6pt, right=6pt, top=5pt, bottom=5pt, #1}

\newtcolorbox{keybox}[1][]{
  breakable, colback=Goldenrod!12, colframe=Goldenrod!70!black,
  boxrule=0.6pt, left=6pt, right=6pt, top=5pt, bottom=5pt, #1}

\newtcolorbox{findbox}[1][]{
  breakable, colback=RedOrange!8, colframe=RedOrange!70!black,
  boxrule=0.6pt, left=6pt, right=6pt, top=5pt, bottom=5pt, #1}

\captionsetup{font=small}
\setlength{\parskip}{4pt}
\linespread{1.15}


\title{\vspace{-1.6cm}پروژه‌ی اول نظریه بازی‌ها\\[4pt]
\large تاب‌آوری راهبردی تحت عدم‌قطعیت:\\
تحلیل بازی‌نظری و پیاده‌سازی در موتور \lr{SEE}}
\author{سهیل سیاح ورگ\\[2pt]
\normalsize شماره‌ی دانشجویی: ۴۰۲۱۱۱۲۳۷\\[4pt]
\normalsize دانشکده‌ی مهندسی کامپیوتر\\[2pt]
\normalsize استاد درس: دکتر حسین قربان}
\date{درس نظریه بازی‌ها --- نیم‌سال دوم سال تحصیلی ۱۴۰۴--۱۴۰۵}

\begin{document}
\maketitle
\vspace{-0.9cm}

\begin{abstract}\noindent
این گزارش چارچوب «تاب‌آوری راهبردی تحت عدم‌قطعیت» را به‌عنوان یک
\emph{جنگ فرسایشی چندمرحله‌ای با اطلاعات ناقص} تحلیل می‌کند و همان تحلیل را
در موتور محاسباتی \lr{SEE} به آزمون می‌گذارد. ادعای مرکزی گزارش این است:
در نمونه‌سازی رسمی این بازی، ارزش خروج ($X\in[-2,46]$) آن‌قدر از جایزه‌ی
دوام‌آوردن ($B\in[87,186]$) کوچک‌تر است که \emph{آستانه‌ی ادامه} به مقدار
بسیار پایینی می‌افتد؛ ما این آستانه را فرمول‌بندی می‌کنیم، مقدار تجربی‌اش را
اندازه می‌گیریم ($\approx 2\%$) و نشان می‌دهیم رفتار عامل آموزش‌دیده دقیقاً با
آن می‌خواند. سه نتیجه‌ی دیگر: (۱) چون قانون هزینه \emph{عمومی} است، پسین روی
مسیر تاب‌آوری حریف تبدیل می‌شود به پسینِ روی یک نوع دوبعدی؛ فیلتر بیزی‌ای که بر
این پایه ساختیم خطای برآورد تاب‌آوری پنهان حریف را $72.9\%$ کاهش می‌دهد.
(۲) قوی‌ترین سیگنال بهترین نیست و تابع سود نسبت به شدت سیگنال \emph{یکنوا
نیست}؛ دلیلش را در قاعده‌ی نوبت‌دهی موتور پیدا کردیم، و توزیع سیگنال عامل
یک \emph{حفره} دقیقاً روی $\sig=0.5$ دارد (‏$29.8\%$ روی $\sig=0.25$، سقوط
به $1.1\%$ روی $\sig=0.5$، و بازگشت به $11.3\%$ روی $\sig=0.75$).
(۳) سه ایراد یا تله در کد ارائه‌شده پیدا کردیم و هر سه را مستند و برطرف
کرده‌ایم: نبود ماژول \lr{see/tournament/runner.py}؛ پیش‌فرض‌های خط فرمان
\lr{train.py} که دقیقاً همان خرابی‌ای را برمی‌گردانند که \lr{DESIGN.md}
می‌گوید رفع شده است؛ و ناهماهنگی استخر حریفِ آموزش با توزیع ارزیابی، که
باعث می‌شد «تسلیم زودهنگام» یک بهترین پاسخ باشد.
(۴) و یک اشتباه روش‌شناختی از خودمان که بعد از اولین بارگذاری روی برد
کلاسی پیدا شد و تصحیحش بیشترین تاثیر را داشت: تمام مدت روی استخر
\emph{سه‌تایی} \lr{scripts/evaluate.py} بهینه می‌کردیم، در حالی که برد
با \emph{شش} حریف مرجع نمره می‌دهد و آن دو معیار در این مسئله در جهت
مخالف هم حرکت می‌کنند. استخر مرجع را بازسازی کردیم (چهار حریف از شش
\emph{دقیقاً} بازتولید می‌شوند) و از دلش دو تصحیح درآمد.
(۵) و یک اشتباه دوم از خودمان، که خواندن دوم برد لو داد: استخر آموزش را
\emph{جایگزین} کرده بودیم به‌جای اینکه \emph{گسترش} بدهیم، و حریف‌های
منفعلِ صبور از آن افتاده بودند. این را هم پیدا و تصحیح کردیم
(بخش~\ref{sec:board2}). امتیاز واقعی برد در سه آپلود
$52.9 \to 56.8 \to \mathbf{58.8}$ شد — بالاتر از هر دو معیار
دستورالعمل، یعنی \lr{TitForTat} با $53.8$ و بنچ‌مارک استاد با $56.3$.
\end{abstract}

\tableofcontents

\vspace{6pt}
\begin{keybox}
\textbf{نقشه‌ی گزارش.} ساختار پیشنهادی بند ۸ دستورالعمل، و اینکه هر بند کجای
این گزارش است:

\vspace{3pt}
\centering\small
\begin{tabular}{@{}rl@{\qquad}rl@{}}
\toprule
بند & جای آن & بند & جای آن\\
\midrule
۱. عنوان پروژه          & صفحه‌ی عنوان        & ۹. اعتبار و انعطاف‌پذیری & بخش~\ref{sec:credflex}\\
۲. چکیده‌ی کوتاه         & چکیده               & ۱۰. ارزش خروج            & بخش~\ref{sec:exitval}\\
۳. بیان مسئله           & بخش~\ref{sec:problem}  & ۱۱. تحلیل تعادلی      & بخش~\ref{sec:equilibrium}\\
۴. خلاصه‌ی مدل           & بخش~\ref{sec:model}    & ۱۲. گزاره‌های تحلیلی   & بخش~\ref{sec:props-stated}\\
۵. تصمیم ادامه یا خروج  & بخش~\ref{sec:contexit} & ۱۳. تفسیر نتایج       & بخش~\ref{sec:interp}\\
۶. هزینه‌ی مؤثر          & بخش~\ref{sec:cost}     & ۱۴. محدودیت‌های مدل    & بخش~\ref{sec:limits}\\
۷. باورها و اطلاعات ناقص& بخش~\ref{sec:beliefs}  & ۱۵. شبیه‌سازی و پیاده‌سازی & بخش~\ref{sec:impl}\\
۸. سیگنال‌دهی پرهزینه    & بخش~\ref{sec:signal}   & ۱۶. جمع‌بندی           & بخش~\ref{sec:concl}\\
\bottomrule
\end{tabular}

\vspace{4pt}
\raggedright\small
بند ۵-۵ دستورالعمل تکمیلی (اتصال گزاره‌ی تحلیلی به داده‌ی عامل) در
بخش~\ref{sec:props} آمده است. تنها تفاوت با ترتیب پیشنهادی این است که بخش
شبیه‌سازی را \emph{قبل} از تفسیر نتایج آورده‌ام، چون تفسیر به عددهای همان بخش
ارجاع می‌دهد.
\end{keybox}

\newpage

% ==================================================================== %
\section{بیان مسئله}\label{sec:problem}
% ==================================================================== %

پرسش اصلی ساده است: \textbf{اگر ادامه‌ی تقابل برای هر دو طرف پرهزینه است، چرا
تقابل ادامه پیدا می‌کند؟}

جواب ساده‌انگارانه («چون یکی از طرفین غیرعقلانی است») در چارچوب این متن رد
می‌شود. متن پیوست نشان می‌دهد ادامه‌ی تقابل می‌تواند برای \emph{هر دو} طرف
عقلانی باشد وقتی چهار چیز هم‌زمان برقرار باشد: طرفین اطلاعات خصوصی درباره‌ی
توان تحمل خود دارند، اقدامات‌شان علاوه بر اثر مادی حامل \emph{اطلاعات} هم
هست، سیگنال معتبر \emph{پرهزینه} است، و ارزش خروج به \emph{تاریخچه} وابسته
است — یعنی می‌شود کاری کرد که خروج دیگر ممکن نباشد.

هدف این گزارش دفاع از هیچ موضع سیاسی نیست؛ هدف این است که ببینیم این چهار
جزء وقتی کنار هم گذاشته می‌شوند چه رفتاری تولید می‌کنند، و آیا آن رفتار در یک
نمونه‌سازی محاسباتی \emph{واقعاً} ظاهر می‌شود یا فقط روی کاغذ است.

% ==================================================================== %
\section{خلاصه‌ی مدل داده‌شده}\label{sec:model}
% ==================================================================== %

بازی به‌صورت
$\Gamma=\langle N,\Theta,p,E,T,H,M,A,\Sigma,D,S,g,G,\varphi,B,X,\mu,\tau,u\rangle$
تعریف شده است. اجزای آن را یک‌به‌یک با شکل رسمی‌شده در
\lr{see/config.py} می‌آوریم.

\paragraph{بازیگران.} $N=\{I,U\}$ (ایران و ایالات متحده).

\paragraph{نوع خصوصی.}
$\theta_i=(\rbar_i,\mbar_i)\in\Theta_i$ که در $t=0$ قرعه‌کشی می‌شود و تا آخر
ثابت می‌ماند: $\rbar_i$ \emph{عزم پایه} و $\mbar_i$ \emph{ظرفیت پایه‌ی مدیریت
هزینه}. پیشین‌ها عمومی و \emph{نامتقارن}‌اند:

\begin{table}[h]\centering
\begin{tabular}{lcc}
\toprule
 & $\rbar$ (عزم) & $\mbar$ (مدیریت هزینه)\\
\midrule
ایران & $\mathrm{Beta}(5,2)$،\ \ $\mathbb{E}\approx0.71$ &
        $0.8+0.9\,\mathrm{Beta}(2,2)$،\ \ $\mathbb{E}\approx1.25$\\
آمریکا & $\mathrm{Beta}(3,3)$،\ \ $\mathbb{E}=0.50$ &
        $0.9+0.9\,\mathrm{Beta}(3,2)$،\ \ $\mathbb{E}\approx1.44$\\
\bottomrule
\end{tabular}
\caption{عدم‌تقارن پیشین‌ها، همان چیزی که بخش ۲ متن می‌گوید: تاب‌آوری ایران از جنس
\emph{جذب فشار} است (عزم بالا) و تاب‌آوری آمریکا از جنس \emph{حفظ فشار}
(مدیریت هزینه‌ی بهتر ولی عزم پراکنده).}
\end{table}

\paragraph{تاب‌آوری و تفاوتش برای دو طرف.}
$e_i^t$ حالت \emph{نهفته}ی تاب‌آوری است، با مقدار اولیه
\[
  e_i^0=100\bigl(0.50+0.35\,\rbar_i+0.15\,\mhat_i\bigr)\in[50,100].
\]
دقت کنید که $\theta_i$ \emph{ریشه‌ی ساختاری} تاب‌آوری است ولی $e_i^t$
\emph{موجودی جاری} آن؛ هر کس تاب‌آوری‌اش به صفر برسد \emph{به‌اجبار} خارج
می‌شود. برای ایران تاب‌آوری یعنی جذب فشار بیرونی بدون بی‌ثباتی درونی؛ برای
آمریکا یعنی نگه‌داشتن فشار بدون خستگی داخلی و کشیدگی راهبردی.

\paragraph{اقدام و سیگنال.}
$a_i^t=(\sig_i^t,d_i^t)$ که
$\sig\in\{0,0.25,0.5,0.75,1\}$ سیگنال عمومی روی یک نردبان پنج‌سطحی است و
$d\in\{C,E\}$ تصمیم ادامه یا خروج. \emph{سیگنال یک اقدام جدا نیست}؛ مؤلفه‌ی
اطلاعاتی همان اقدام است.

\paragraph{قاعده‌ی نوبت.}
$t=0$ هم‌زمان؛ اگر در $t-1$ \emph{دقیقاً یکی} از طرفین $\sig\ge0.75$ بازی کرده
باشد، مرحله‌ی بعد فقط طرف \emph{مقابل} حرکت می‌کند و تشدیدکننده مجبور به
\lr{HOLD} است — یعنی موضعش عمومی باقی می‌ماند و \emph{نمی‌تواند آن مرحله خارج
شود}. در غیر این صورت هم‌زمان.

\paragraph{هزینه‌ی مؤثر.} دقیقاً همان‌طور که در متن آمده:
\[
  g_i^t=\frac{c_i(t,h^t)+a_i^{\mathrm{aud}}(t,h^t)+r_i(t,h^t)}{m_i^t},
\]
با تصریح
\[
  c_i=c_{0i}+\kappa_i\tilde\sig_j,\qquad
  a_i^{\mathrm{aud}}=\alpha_i K_i^t,\qquad
  r_i=\eta\,\tilde\sig_i\tilde\sig_j,\qquad
  m_i^t=\mbar_i\bigl(0.35+0.65\,\ehat_i^t\bigr),
\]
که $K_i^{t+1}=0.75K_i^t+\sig_i^t$ \emph{انبار تعهد عمومی} است
(‏$K\le4$) و کل عبارت در شیب تشدید $\min\bigl(1,\tfrac{t+1}{4}\bigr)$ ضرب
می‌شود. قانون گذار:
$e_i^{t+1}=e_i^t-g_i^t+\xi^t$ با $\xi\sim\mathcal N(0,0.5^2)$.

\paragraph{هزینه‌ی سیگنال‌دهی.}
\[
  \varphi_i(\sig)=\frac{1.2\,\sig^2}
  {(0.4+0.6\rbar_i)\,(0.4+0.6\min(\ehat_i,1))\,(0.5+0.5\mhat_i)} .
\]

\paragraph{دو ارزش پایانی.}
\[
  B_i=175\bigl(0.50+0.30\,\ehat_i+0.20\,\rbar_i\bigr),\qquad
  X_i(t,h^t)=x_{0i}-5K_i^t-8\tilde\sig_j+20M^t .
\]
$M^t\in\{0,1\}$ پنجره‌ی میانجی‌گری است: بسته $\to$ باز با احتمال $0.35$ اگر
موضع مشترک آرام باشد ($\tilde\sig_I+\tilde\sig_U\le0.5$) وگرنه $0.05$؛ باز
$\to$ بسته با احتمال $0.25$.

\paragraph{باورها و عایدی.}
$\mu_j^t(\theta_i,e_i^t\mid h^t)$ پسین $j$ درباره‌ی $i$ است. با
$\tau_i=\inf\{t:d_i^t=E\}$:
\[
  u_i=
  \begin{cases}
    B_i(e_i^{\tau_j},\theta_i)-G_i(\tau_j)-\sum_{s\le\tau_j}\varphi_i(\sig_i^s),
      & \tau_i>\tau_j,\\[4pt]
    X_i(\tau_i,h^{\tau_i})-G_i(\tau_i)-\sum_{s\le\tau_i}\varphi_i(\sig_i^s),
      & \tau_i\le\tau_j .
  \end{cases}
\]

\begin{keybox}
\textbf{ده جزئی که بند ۳ دستورالعمل خواسته، و جای هرکدام.}

\vspace{3pt}
\centering\small
\begin{tabular}{@{}rl@{\qquad}rl@{}}
\toprule
جزء & کجاست & جزء & کجاست\\
\midrule
بازیگران اصلی & $N=\{I,U\}$، بالا & هزینه‌ی مؤثر & $g_i^t$، بالا\\
نوع خصوصی & $\theta_i=(\rbar_i,\mbar_i)$، بالا & هزینه‌ی سیگنال‌دهی & $\varphi_i(\sig)$، بالا\\
تاب‌آوری و تفاوتش & $e_i^t$، بالا & ارزش دوام تا عقب‌نشینی حریف & $B_i$، بالا\\
اقدام‌ها و سیگنال‌ها & $a_i^t=(\sig_i^t,d_i^t)$ & ارزش خروج/مصالحه & $X_i(t,h^t)$، بالا\\
تصمیم ادامه یا خروج & $d_i^t\in\{C,E\}$ & باورها و به‌روزرسانی & $\mu_j^t$، و بخش~\ref{sec:beliefs}\\
\bottomrule
\end{tabular}

\vspace{4pt}
\raggedright\small
این بخش عمداً کوتاه است و نه مرور ادبیات است و نه مدل تازه‌ای می‌سازد؛ فقط
همان مدل متن را با شکل عددی‌شده‌ی \lr{see/config.py} کنار هم می‌گذارد.
\end{keybox}

\begin{keybox}
\textbf{یک رابطه‌ی مفید که بارها استفاده می‌کنیم.} چون
$e_i^{t+1}=e_i^t-g_i^t+\xi$، داریم $G_i(t)\approx e_i^0-e_i^t$. یعنی
«هزینه‌ی مؤثر انباشته» و «تاب‌آوری خرج‌شده» \emph{یک چیزند}. پس
\[
  u_i^{\text{دوام}}\ \approx\ B_i(\ehat_i)-\bigl(e_i^0-e_i\bigr)
  -\textstyle\sum\varphi_i .
\]
\end{keybox}

% ==================================================================== %
\section{تحلیل تصمیم ادامه یا خروج}\label{sec:contexit}
% ==================================================================== %

در هر مرحله بازیکن $i$ بین ادامه و خروج انتخاب می‌کند. نکته‌ی اول این است که
هزینه‌های گذشته \emph{ریخته}‌اند: هر دو شاخه یک $-G_i(t)-\sum_{s<t}\varphi_i$
مشترک دارند. پس مقایسه کاملاً \emph{آینده‌نگر} است:
\[
  \underbrace{q_i^t\,\mathbb{E}[B_i]+(1-q_i^t)\,\mathbb{E}[X_i^{\text{بعدی}}]
  -\mathbb{E}\Bigl[\textstyle\sum_{s\ge t}(g_i^s+\varphi_i^s)\Bigr]}
  _{\text{ارزش ادامه } V_i^C}
  \ \ \ge\ \ \underbrace{X_i(t,h^t)}_{\text{ارزش خروج امروز}},
\]
که $q_i^t=\Pr(\tau_j<\tau_i\mid h^t,\mu_i^t)$ احتمال \emph{باورمحور} برنده
شدن مسابقه‌ی فرسایش است.

\subsection{دروازه‌ی ادامه}

نسخه‌ی یک‌مرحله‌ای این شرط، فرمول‌بندی مفیدی می‌دهد. اگر یک مرحله‌ی دیگر بمانیم:

\begin{proposition}[دروازه‌ی ادامه]\label{pr:gate}
ادامه دادن به‌صرفه است اگر و تنها اگر
\[
  q_i^t\ \ge\ q^\ast_i(t,h^t)
  \ :=\ \frac{g_i^t+\varphi_i(\sig_i^t)-\Delta X_i}{B_i-X_i},
  \qquad
  \Delta X_i:=\mathbb{E}\bigl[X_i(t{+}1,h^{t+1})\bigr]-X_i(t,h^t).
\]
\end{proposition}

سه چیز در این کسر هست و هر سه معنی دارند:
\begin{itemize}[leftmargin=1.4em,itemsep=2pt]
  \item \textbf{صورت} هزینه‌ی یک مرحله ماندن است: سوخت جاری $g$، بهای
        سیگنال $\varphi$، و $-\Delta X_i$ که \emph{فرسایش گزینه‌ی خروج} است.
        این جمله‌ی سوم همان چیزی است که در تحلیل‌های ساده جا می‌افتد: ماندن
        فقط پول خرج نمی‌کند، \emph{در خروج را هم تنگ‌تر می‌کند}.
  \item \textbf{مخرج} فاصله‌ی جایزه تا گزینه‌ی بیرونی است.
  \item اگر $q$ از این کسر بیشتر باشد، ادامه عقلانی است — \emph{حتی اگر
        ادامه پرهزینه باشد}. این پاسخ مستقیم به پرسش «چرا ممکن است یک بازیگر
        ادامه دهد حتی اگر ادامه برایش پرهزینه باشد؟».
\end{itemize}

\subsection{عدد این دروازه در جهان \lr{SEE}}

حالا بازه‌های واقعی را بگذاریم:
\[
  B\in[87.5,\,185.5],\qquad X_I\in[-2,\,46],\qquad X_U\in[0,\,48],
\]
و در یک موضع آرام ($\sig=0$، $K=0$، $e=e^0$) هزینه‌ی جاری فقط
$g\approx1.12$ (ایران) و $g\approx1.33$ (آمریکا) در هر مرحله است. پس با
$B\approx150$ و $X\approx26$:
\[
  q^\ast\approx\frac{1.12}{150-26}\approx0.009 .
\]

\begin{findbox}
\textbf{یافته‌ی کلیدی ۱.} در این نمونه‌سازی، دروازه‌ی ادامه \emph{بسیار پایین}
است. بازیکن باید ادامه دهد مگر آنکه تقریباً مطمئن باشد بازنده‌ی مسابقه‌ی
فرسایش است. به بیان دیگر: \textbf{خروج داوطلبانه تقریباً هیچ‌وقت بهینه نیست؛
خروجِ عقلانی همان خروج \emph{اجباری} است.} این نتیجه بعداً هم به‌صورت آماری
روی رفتار عامل (میانه‌ی $q^\ast\approx2\%$) و هم به‌صورت مسابقه‌ای
(بخش~\ref{sec:sweep}) تایید می‌شود.
\end{findbox}

% ==================================================================== %
\section{تحلیل هزینه‌ی مؤثر و ظرفیت مدیریت هزینه}\label{sec:cost}
% ==================================================================== %

هزینه‌ی ادامه فقط هزینه‌ی خام نیست. سه جزء در صورت کسر است — مادی
($c_i$)، داخلی/مخاطبان ($a^{\mathrm{aud}}_i$)، و ریسک تشدید ($r_i$) — و همه
تقسیم بر \emph{ظرفیت مدیریت هزینه} $m_i^t$ می‌شوند.

\paragraph{افزایش هزینه‌ی خام چه می‌کند؟}
اثر مستقیم و خطی: $g$ متناسب با صورت بالا می‌رود. ولی توجه کنید که
$c_i=c_{0i}+\kappa_i\tilde\sig_j$ به موضع \emph{طرف مقابل} بسته است. یعنی
فشار مادی چیزی است که حریف بر شما تحمیل می‌کند، نه چیزی که خودتان انتخاب
می‌کنید.

\paragraph{افزایش ظرفیت مدیریت چه می‌کند؟}
اثرش \emph{ضربی} و در نتیجه قوی‌تر است، و یک بازخورد هم دارد:
$m_i^t=\mbar_i(0.35+0.65\ehat_i^t)$ خودش با فرسایش تاب‌آوری کم می‌شود. پس
همان فشار خام وقتی از قبل فرسوده‌اید بیشتر می‌سوزاند — این همان
\emph{شکنندگی فزاینده} است.

\subsection{زمان فروپاشی: یک فرمول بسته}

این بازخورد را می‌شود دقیق حل کرد. با موضع ثابت، صورت کسر ($R$) ثابت است و با
$x=e/e^0$:
\[
  \frac{dx}{dt}=\frac{-R}{e^0\,\mbar\,(0.35+0.65x)}
  \ \Longrightarrow\
  \int_0^x(0.35+0.65u)\,du=\frac{R}{e^0\mbar}\,T ,
\]
و در نتیجه

\begin{proposition}[زمان تا خروج اجباری]\label{pr:collapse}
\[
  \boxed{\;T_{\text{فروپاشی}}
  =\frac{e^0\,\mbar\,\bigl(0.35\,x+0.325\,x^2\bigr)}{R}\;}
  \qquad (x=\ehat).
\]
\end{proposition}

سه چیز از این فرمول می‌خوانیم:
\begin{enumerate}[leftmargin=1.6em,itemsep=2pt]
  \item $T$ نسبت به $\mbar$ \emph{خطی} است ولی نسبت به فشار خام $R$
        \emph{وارون}. یعنی دو برابر کردن ظرفیت مدیریت همان‌قدر عمر می‌خرد که
        نصف کردن فشار — ولی معمولاً اولی در اختیار خودتان است و دومی نیست.
  \item $T$ نسبت به تاب‌آوری باقی‌مانده \emph{محدب} است ($0.35x+0.325x^2$):
        از دست دادن $10\%$ اول تاب‌آوری خیلی کم‌خطرتر از $10\%$ آخر است.
        این ریاضیاتِ همان «شکنندگی فزاینده» است.
  \item این فرمول در بخش پیاده‌سازی به‌عنوان آماره‌ی تصمیم استفاده می‌شود
        (بخش~\ref{sec:impl})، پس تحلیل مستقیماً به کد وصل می‌شود.
\end{enumerate}

\paragraph{چرا فشار خارجی لزوماً تاب‌آوری را کم نمی‌کند؟}
چون فشار از کانال $R/m$ وارد می‌شود. اگر بازیگری بتواند فشار بیرونی را به
هزینه‌ی \emph{سیاسی قابل مدیریت} تبدیل کند (یعنی $m$ را بالا نگه دارد)،
افزایش $R$ می‌تواند تقریباً خنثی شود.

\paragraph{چرا ضعف در مدیریت داخلی تاب‌آوری راهبردی را کم می‌کند؟}
دو کانال: اول، $m$ کوچک‌تر یعنی هر واحد فشار بیشتر می‌سوزاند. دوم و مهم‌تر،
جمله‌ی $a^{\mathrm{aud}}_i=\alpha_i K_i$ نشان می‌دهد بخشی از هزینه اصلاً از
بیرون نمی‌آید — \emph{تعهدات عمومی خودِ بازیگر} است. یعنی می‌شود بدون هیچ
اقدامی از طرف مقابل، صرفاً با بلندپروازی در حرف، هزینه‌ی خود را بالا برد.

% ==================================================================== %
\section{تحلیل باورها و اطلاعات ناقص}\label{sec:beliefs}
% ==================================================================== %

\paragraph{هر بازیکن درباره‌ی چه چیزی عدم‌قطعیت دارد؟}
درباره‌ی $\theta_j=(\rbar_j,\mbar_j)$ و در نتیجه درباره‌ی $e_j^t$. عامل فقط
نوع \emph{خودش}، تاب‌آوری \emph{خودش} و کل تاریخچه‌ی \emph{عمومی} را می‌بیند.

\subsection{مشاهده‌ی اصلی: مسئله‌ی فیلتر دوبعدی است}

\begin{keybox}
\textbf{ایده‌ای که کل بخش پیاده‌سازی روی آن سوار است.}
حالت $e_j^t$ پنهان است، ولی \emph{قانونی} که آن را حرکت می‌دهد عمومی است:
\[
  g_j^t=\mathrm{ramp}(t)\cdot
  \frac{c_{0j}+\kappa_j\sig_i+\alpha_j K_j+\eta\sig_i\sig_j}
       {\mbar_j\,(0.35+0.65\,e_j/e_j^0)} .
\]
هر چیزی در سمت راست، جز $(\mbar_j,e_j^0)$، یا یک ثابت منتشرشده است یا جزئی
از تاریخچه‌ی عمومی. پس \textbf{به شرط یک نوع فرضی $\theta_j$، کل مسیر
تاب‌آوری حریف تعیین‌شده است} (تا حد نویز‌ی کوچک $\xi$).
یعنی پسین روی \emph{مسیر}، تبدیل می‌شود به پسین روی یک \emph{نوع دوبعدی}.
\end{keybox}

\subsection{چگونه مشاهده‌ی گذشته باورها را عوض می‌کند}

سه کانال به‌روزرسانی داریم:

\paragraph{(الف) بقا — و چرا «هنوز خارج نشدن» خودش خبر است.}
اگر $e_j\le0$ شود، حریف \emph{به‌اجبار} خارج می‌شود. پس تا وقتی حریف سر جایش
است، هر نوعی که شبیه‌سازی‌اش تا حالا به صفر رسیده باشد \emph{ناممکن} است و
وزنش صفر می‌شود. این کانال \emph{دقیق} است و هیچ مدلی از استراتژی حریف لازم
ندارد — فقط قانون هزینه. این دقیقاً همان «اطلاعاتی که در نفس این واقعیت
نهفته است که طرف مقابل هنوز خارج نشده» است.

\paragraph{(ب) سیگنال پرهزینه.}
چون $\varphi$ برای نوع‌های قوی ارزان‌تر است، موضع بلندِ \emph{انتخاب‌شده}
شاهدی به نفع قوی بودن است. ما این را با درست‌نمایی پاسخ کوانتال
$\propto\exp(-\beta\varphi_j(\sig_j;\theta_j,\ehat_j))$ وارد می‌کنیم، و فقط
در مرحله‌هایی که حریف \emph{واقعاً انتخاب کرده} (نه \lr{HOLD} اجباری).

\paragraph{(ج) پیشین منتشرشده}، که نامتقارن است.

\subsection{یک مثال ساده از به‌روزرسانی باور}

فرض کنید هر دو طرف تا مرحله‌ی $t=6$ موضع آرام ($\sig=0$) داشته‌اند و آمریکا
هنوز خارج نشده. با $\sig=0$، فشار خام روی آمریکا
$R_U=c_{0U}=1.8$ است و طبق گزاره~\ref{pr:collapse} یک نوع با
$\mbar_U$ و $e_U^0$ عمری برابر
$T=e_U^0\mbar_U(0.675)/1.8$ دارد. برای اینکه تا $t=6$ دوام بیاورد باید
$T>6$، یعنی
\[
  e_U^0\,\mbar_U\ >\ \frac{6\times1.8}{0.675}=16 .
\]
چون $e_U^0\ge50$ و $\mbar_U\ge0.9$، این شرط در $t=6$ هنوز هیچ نوعی را حذف
نمی‌کند. اما اگر آمریکا هفت مرحله زیر موضع $\sig_I=1$ دوام بیاورد
($R_U=1.8+2.5=4.3$)، شرط به $e_U^0\mbar_U>44.6$ سفت می‌شود و بخشی از
پایین‌ترین نوع‌ها حذف می‌شوند. \textbf{یعنی فشار بیشتر، اطلاعات بیشتری هم
تولید می‌کند} — یک انگیزه‌ی اطلاعاتی برای تشدید که در متن اصلی صریح نیست.

\paragraph{چرا عقب‌نشینی نکردن همیشه معنای قدرت واقعی نیست؟}
چون دو نوع مختلف می‌توانند رفتار یکسان تولید کنند. یک نوع ضعیف با $\ehat$
بالا (تازه‌نفس ولی کم‌عزم) دقیقاً مثل یک نوع قوی رفتار می‌کند تا وقتی فشار
کافی جمع نشده. تفکیک فقط وقتی رخ می‌دهد که هزینه‌ی انباشته به‌قدری بالا برود
که تقلید برای نوع ضعیف گران شود.

\paragraph{چگونه باور اشتباه تقابل را طولانی می‌کند؟}
اگر $\mu_i$ به اشتباه $q_i$ را بالا برآورد کند، طبق
گزاره~\ref{pr:gate} بازیکن ادامه می‌دهد. اگر \emph{هر دو} طرف چنین خطایی
داشته باشند، هر دو ادامه می‌دهند و تقابل تا وقتی طول می‌کشد که واقعیت
(فرسایش تاب‌آوری) باور را اصلاح کند. چون دروازه‌ی ادامه پایین است
(‏$q^\ast\approx0.01$)، حتی خطای کوچک در باور هم می‌تواند تقابل را طولانی
کند.

% ==================================================================== %
\section{تحلیل سیگنال‌دهی پرهزینه}\label{sec:signal}
% ==================================================================== %

\paragraph{چرا سیگنال بی‌هزینه ممکن است معتبر نباشد؟}
اگر $\varphi\equiv0$ بود، هر نوعی می‌توانست $\sig=1$ بفرستد بدون اینکه چیزی
از دست بدهد. آن‌وقت سیگنال هیچ اطلاعاتی حمل نمی‌کرد و تنها تعادل، تعادل
\emph{تجمیعی} با بلندترین سیگنال بود. اعتبار از این می‌آید که \emph{هزینه‌ی
سیگنال برای نوع‌های مختلف فرق می‌کند}، نه از خودِ پیام.

\paragraph{چرا سیگنال قوی برای نوع ضعیف پرهزینه‌تر است؟}
مستقیماً از شکل $\varphi$: مخرج در $\rbar$، $\ehat$ و $\mhat$ صعودی است، پس
$\partial\varphi/\partial\rbar<0$ و $\partial\varphi/\partial\ehat<0$.
عدد این تفاوت بزرگ است:

\begin{table}[h]\centering
\begin{tabular}{lc}
\toprule
نوع & $\varphi(\sig=1)$\\
\midrule
قوی و تازه‌نفس ($\rbar=0.95$، $\ehat=1$، $\mhat=0.95$) & $1.27$\\
ضعیف و فرسوده ($\rbar=0.10$، $\ehat=0.2$، $\mhat=0.10$) & $9.12$\\
\midrule
\textbf{نسبت} & $\mathbf{7.2\times}$\\
\bottomrule
\end{tabular}
\caption{شکاف هزینه‌ی سیگنال بین دو سر طیف نوع. برای مقایسه، هزینه‌ی جاری
یک مرحله در موضع آرام حدود $1.1$ است — پس این شکاف در مقیاس بازی
\emph{بزرگ} است.}
\end{table}

\paragraph{چگونه شرط \lr{single-crossing} سیگنال را اطلاعاتی می‌کند؟}
شرط لازم این است که \emph{هزینه‌ی نهایی} افزایش شدت سیگنال برای نوع قوی
کمتر باشد:
\[
  \frac{\partial^2\varphi_i}{\partial\sig\,\partial\ehat_i}<0,
  \qquad
  \frac{\partial^2\varphi_i}{\partial\sig\,\partial\rbar_i}<0 .
\]
برای شکل رسمی‌شده، $\partial\varphi/\partial\sig=2.4\sig/D$ که در آن $D$ در
$\rbar$ صعودی است، پس شرط برقرار است. عدداً در $\sig=0.5$:

\begin{table}[h]\centering
\begin{tabular}{lccc}
\toprule
$\rbar$ & $0.2$ & $0.5$ & $0.8$\\
\midrule
$\partial\varphi/\partial\sig$ & $3.08$ & $2.29$ & $1.82$\\
\bottomrule
\end{tabular}
\caption{هزینه‌ی نهایی شدت سیگنال با عزم پایه \emph{نزولی} است — همان شرط
تک‌تقاطعی (اسپنس-میرلیز). آزمون
\lr{tests/test\_env.py::test\_single\_crossing\_of\_signaling\_cost}
همین را عددی تایید می‌کند.}
\end{table}

معنای این شرط: چون نوع ضعیف برای بالا بردن سیگنال \emph{بیشتر} می‌پردازد،
مجموعه‌ای از سیگنال‌ها هست که نوع قوی حاضر است بفرستد و نوع ضعیف نیست. همان
جا تعادل \emph{تفکیکی} ممکن می‌شود.

\paragraph{تفاوت سیگنال واقعی و بلوف.}
بلوف یعنی نوع ضعیف $\sig$ بالا می‌فرستد. در این مدل بلوف \emph{ممکن} است ولی
\emph{دوگانه پرهزینه}: هم $\varphi$ بیشتری می‌پردازد (که تاب‌آوری‌اش را سریع‌تر
می‌سوزاند)، و هم $K$ را بالا می‌برد که هم $a^{\mathrm{aud}}$ را زیاد و هم
$X$ را کم می‌کند. پس بلوف در کوتاه‌مدت باور حریف را جابه‌جا می‌کند ولی در
میان‌مدت خودش را لو می‌دهد: نوع ضعیفِ بلوف‌زننده \emph{زودتر} فرو می‌پاشد.
این دقیقاً همان چیزی است که کانال (الف) در بخش~\ref{sec:beliefs} می‌گیرد.

% ==================================================================== %
\section{تحلیل تعارض اعتبار و انعطاف‌پذیری}\label{sec:credflex}
% ==================================================================== %

سیگنال قوی‌تر اعتبار می‌سازد، ولی هم‌زمان امکان خروج را تنگ می‌کند. در این
نمونه‌سازی، این تعارض از \emph{سه} کانال مشخص می‌گذرد:

\begin{enumerate}[leftmargin=1.6em,itemsep=3pt]
  \item \textbf{انبار تعهد.} $K_i^{t+1}=0.75K_i^t+\sig_i^t$. یک موضع
        \emph{پایدار} $\sig$ در بلندمدت به $K=4\sig$ می‌رسد، و چون
        $X_i=x_{0i}-5K_i-\dots$، یک موضع پایدار $\sig$ در حالت پایدار
        $20\sig$ واحد از ارزش خروج می‌خورد. با $x_0=26$ این یعنی
        \textbf{یک موضع پایدار $\sig=1$ تقریباً کل گزینه‌ی خروج را نابود
        می‌کند}.
  \item \textbf{تحقیر.} جمله‌ی $-8\tilde\sig_j$: خروج زیر موضع بلند طرف
        مقابل گران‌تر است.
  \item \textbf{قفل نوبت.} اگر $\sig_i\ge0.75$ و حریف نه، مرحله‌ی بعد فقط
        حریف حرکت می‌کند و شما \emph{به‌اجبار} \lr{HOLD} می‌کنید — یعنی
        حتی اگر بخواهید، \emph{نمی‌توانید} آن مرحله خارج شوید. این
        صریح‌ترین شکل «گیر افتادن در تعهد خود» است.
\end{enumerate}
و یک کانال چهارم غیرمستقیم: موضع بلند احتمال باز شدن پنجره‌ی میانجی‌گری را از
$0.35$ به $0.05$ می‌اندازد، و هر پنجره‌ی باز $+20$ به $X$ اضافه می‌کند.

\subsection{اما تخفیفی که موتور به تشدید می‌دهد}

اینجا به چیزی رسیدیم که در تحلیل روی کاغذ دیده نمی‌شود و فقط با خواندن دقیق
پیاده‌سازی پیدا شد.

\begin{findbox}
\textbf{یافته‌ی کلیدی ۲ — یکنوا نبودنِ سود نسبت به شدت سیگنال.}
موتور $\varphi$ را \emph{فقط وقتی یک سیگنال تازه تولید می‌شود} دریافت می‌کند؛
در حالت \lr{HOLD} موضع حفظ می‌شود ولی \textbf{هیچ هزینه‌ی سیگنال تازه‌ای گرفته
نمی‌شود}. از طرف دیگر، عبور از آستانه‌ی $\sig\ge0.75$ دقیقاً همان چیزی است که
شما را مجبور به \lr{HOLD} می‌کند.

نتیجه: یک موضع پایدار $\sig=0.75$ تقریباً \emph{یک مرحله در میان} صورت‌حساب
سیگنال می‌دهد، در حالی که موضع $\sig=0.5$ \emph{هر مرحله} می‌پردازد. با
$\varphi\propto\sig^2$:
\[
  \tfrac12\varphi(0.75)=\tfrac12(0.5625)\kappa=0.281\kappa
  \quad\text{در برابر}\quad
  \varphi(0.50)=0.25\kappa ,
\]
یعنی هزینه‌ی سیگنال \emph{عملاً برابر}، ولی فشار وارد بر حریف
($\kappa\tilde\sig$) و ریسک تشدید ($\eta\tilde\sig_i\tilde\sig_j$) با
$\sig=0.75$ به‌مراتب بیشتر.

پس تابع سود نسبت به $\sig$ \textbf{یکنوا نیست}: $\sig=0.5$ بدترین گزینه است
(هزینه‌ی هر مرحله بدون مزیت قفل)، و $\sig=0.75$ بهترین نقطه است.
\end{findbox}

این یک \emph{انحراف} نمونه‌سازی از قصد متن است: متن می‌خواهد سیگنال قوی‌تر
گران‌تر و خودمحدودکننده باشد، ولی موتور با معافیت $\varphi$ در \lr{HOLD}
بخشی از آن گرانی را پس می‌دهد. ما این را به‌عنوان یک محدودیت گزارش می‌کنیم
(بخش~\ref{sec:limits}) و در بخش تجربی نشان می‌دهیم عامل آموزش‌دیده دقیقاً همین
را کشف می‌کند.

\paragraph{پس چرا قوی‌ترین سیگنال بهترین نیست؟}
چون $\sig=1$ از یک طرف همان مزیت قفل $\sig=0.75$ را دارد ولی از طرف دیگر
$\varphi\propto\sig^2$ را $\tfrac{1}{0.5625}=1.78$ برابر می‌کند، $K$ و در
نتیجه هزینه‌ی مخاطبان را بیشتر بالا می‌برد، و ارزش خروج را بیشتر می‌خورد. در
سنجش تجربی (بخش~\ref{sec:sweep}) $\sig=0.75$ حدود $13$ واحد بهتر از
$\sig=1.0$ است.

\paragraph{چگونه هم معتبر باشیم و هم مسیر خروج را نبندیم؟}
پاسخ مدل: \emph{موضع بلند را پایدار نکنید}. چون $K$ با ضریب $0.75$ فراموش
می‌کند، یک تشدید کوتاه $K$ را زیاد بالا نمی‌برد ولی همان یک مرحله برای
فرستادن سیگنال و گرفتن نوبت پاسخ کافی است. اعتبار از \emph{رخداد} تشدید
می‌آید، فرسایش خروج از \emph{انباشت} آن.

% ==================================================================== %
\section{تحلیل ارزش خروج}\label{sec:exitval}
% ==================================================================== %

$X_i(t,h^t)=x_{0i}-5K_i^t-8\tilde\sig_j+20M^t$. خروج در این مدل مساوی شکست
نیست؛ می‌تواند مصالحه، آتش‌بس یا تغییر راهبرد باشد. سه جمله سه مکانیزم‌اند:

\paragraph{چرا ارزش خروج به تاریخچه وابسته است؟}
چون $K_i^t$ خلاصه‌ی تمام تعهدات عمومی گذشته‌ی خودتان است و $\tilde\sig_j$
موضع فعلی حریف. یعنی همان تاریخچه‌ای که اعتبار ساخته، در خروج را تنگ کرده.

\paragraph{در کدام تاریخچه‌ها خروج آسان‌تر است؟}
وقتی $K$ پایین است (تعهد عمومی کم داده‌اید)، حریف موضع آرام دارد، و پنجره‌ی
میانجی‌گری باز است. توجه کنید هر سه شرط با \emph{خویشتن‌داری متقابل} هم‌جهت‌اند
— و پنجره‌ی میانجی‌گری هم فقط زیر خویشتن‌داری متقابل با احتمال بالا باز
می‌شود ($0.35$ در برابر $0.05$).

\paragraph{چگونه هزینه‌های داخلی خروج را دشوار می‌کنند؟}
جمله‌ی $-5K_i$. این تنها جایی است که \emph{انتخاب گذشته‌ی خودِ بازیگر}
مستقیماً گزینه‌ی امروزش را می‌خورد. با $K\le4$، این جمله تا $-20$ می‌رود، در
حالی که $x_0$ فقط $26$ است.

\paragraph{نقش میانجی و زبان \lr{face-saving}.}
جمله‌ی $+20M^t$ بزرگ‌ترین تک‌جمله‌ی این معادله است — تقریباً به اندازه‌ی کل
$x_0$. یعنی در این مدل، \emph{ساختن} یک مسیر خروج آبرومندانه می‌تواند به‌تنهایی
ارزش خروج را دو برابر کند. این شکل ریاضی همان حرف است که «خروج باید ساخته شود، نه اینکه از
قبل موجود باشد» است.

\paragraph{چرا ممکن است بازیگر ادامه دهد نه به امید پیروزی؟}
طبق گزاره~\ref{pr:gate}، وقتی $X_i$ به‌خاطر $K$ بالا و $M=0$ نزدیک صفر
می‌شود، مخرج $B-X$ بزرگ می‌شود و $q^\ast$ باز هم کوچک‌تر. یعنی
\textbf{هرچه خروج بی‌آبروتر باشد، ادامه دادن عقلانی‌تر می‌شود} — حتی با
احتمال پیروزی ناچیز. این حلقه‌ی خودتقویت‌کننده مکانیزم اصلی طولانی شدن
تقابل در این مدل است.

% ==================================================================== %
\section{تحلیل تعادلی}\label{sec:equilibrium}
% ==================================================================== %

مفهوم تعادلی متناسب با این بازی \emph{تعادل بیزی کامل} (\lr{PBE}) است:
استراتژی‌ها باید در هر زیربازی \emph{عقلانیت ترتیبی} داشته باشند و باورها
هرجا ممکن است با قاعده‌ی بیز به‌روز شوند.

\subsection{ساختار آستانه‌ای}

\begin{proposition}[تعادل از نوع آستانه‌ای]\label{pr:cutoff}
تاریخچه‌ی $h^t$ و باور $\mu_i^t$ را ثابت بگیرید و
$\Delta_i(\theta_i,e_i):=V_i^C(\theta_i,e_i,h^t,\mu_i^t)-X_i(t,h^t)$
را مازاد ادامه بنامید. آنگاه $\Delta_i$ نسبت به $e_i$، $\rbar_i$ و $\mbar_i$
اکیداً صعودی است. در نتیجه استراتژی بهینه \emph{آستانه‌ای} است: بازیکن ادامه
می‌دهد اگر و تنها اگر شاخص قدرتش از یک آستانه‌ی وابسته به تاریخچه و باور
بیشتر باشد.
\end{proposition}

\begin{proof}[طرح اثبات]
چهار مشتق را جمع می‌زنیم. $B_i=175(0.5+0.3\ehat_i+0.2\rbar_i)$ در $e_i$ و
$\rbar_i$ صعودی است. هزینه‌ی جاری $g_i=R/(\mbar_i(0.35+0.65\ehat_i))$ در
$\mbar_i$ و $e_i$ نزولی است. هزینه‌ی سیگنال $\varphi_i$ در هر سه‌ی
$\rbar_i,\ehat_i,\mhat_i$ نزولی است. و $X_i(t,h^t)$ اصلاً به $\theta_i$
بستگی ندارد. پس هر چهار جمله در یک جهت کار می‌کنند و $\Delta_i$ صعودی است؛
یکنوایی مازاد نسبت به نوع، شرط کافی برای آستانه‌ای بودن قاعده‌ی توقف است.
\end{proof}

\paragraph{پاسخ به شش پرسش بخش ۴-۷ دستورالعمل.}
\begin{enumerate}[leftmargin=1.7em,itemsep=3pt]
  \item \textbf{استراتژی نوع با تاب‌آوری بالا:} ادامه بده و در مراحل اولیه
        تشدید کن. دلیل: $\varphi$ برایت ارزان است (پس سیگنال کم‌هزینه است)،
        $B$ برایت بزرگ است، و طبق گزاره~\ref{pr:collapse} عمرت بلند است.
        تشدید هم‌زمان دو کار می‌کند: باور حریف را جابه‌جا می‌کند و
        $T_{\text{فروپاشی}}$ او را کوتاه می‌کند.
  \item \textbf{استراتژی نوع با تاب‌آوری پایین:} $K$ را پایین نگه دار، موضع
        آرام بگیر تا پنجره‌ی میانجی‌گری باز شود، و \emph{آن‌وقت} خارج شو.
        چون $X$ با $M=1$ حدود دو برابر می‌شود، این تنها راه خروج کم‌ضرر است.
        نکته اینکه نوع ضعیف \emph{نباید} زود خارج شود — طبق
        گزاره~\ref{pr:gate} حتی او هم باید صبر کند تا یا میانجی برسد یا
        اجباراً خارج شود.
  \item \textbf{اثر باور حریف:} از طریق $q_i^t$ در گزاره~\ref{pr:gate}. اگر
        حریف باور کند شما قوی‌اید، $q_j$ او پایین می‌آید و زودتر به آستانه‌ی
        خروجش می‌رسد.
  \item \textbf{چه وقت ادامه برای هر دو پایدار می‌ماند؟} وقتی
        $q_I^t\ge q_I^\ast$ و $q_U^t\ge q_U^\ast$ هم‌زمان. چون
        $q_I+q_U\approx1$ ولی $q^\ast\approx0.01$، این وضعیت \emph{تقریباً
        همیشه} برقرار است. \textbf{این پاسخ رسمی به پرسش اصلی پروژه است:}
        تقابل ادامه پیدا می‌کند چون دو آستانه‌ی خیلی پایین می‌توانند
        هم‌زمان ارضا شوند.
  \item \textbf{چه وقت یکی مایل به خروج می‌شود؟} وقتی $q$ او زیر آستانه
        برود (نزدیک فروپاشی) یا آستانه‌اش بالا برود — یعنی $X$ ناگهان بالا
        بپرد، که دقیقاً همان باز شدن پنجره‌ی میانجی‌گری است.
  \item \textbf{تفکیکی یا تجمیعی؟} هر دو ممکن است. شرط تک‌تقاطعی امکان
        تعادل \emph{تفکیکی} را می‌دهد، ولی چون نردبان سیگنال \emph{گسسته}
        است (پنج سطح) و آستانه‌ی $0.75$ یک ناپیوستگی دارد، آنچه عملاً
        دیدیم یک تعادل \emph{نیمه‌تجمیعی} است: توده‌ی احتمال روی دو سطح
        $\{0,\,0.75\}$ جمع می‌شود و سطوح میانی تقریباً استفاده نمی‌شوند
        (بخش~\ref{sec:props}: $99.2\%$ مراحل روی همین دو سطح، و
        $\sig=0.5$ فقط $0.2\%$).
\end{enumerate}

\begin{remark}
آنچه با خودبازی به دست می‌آید یک \emph{تعادل تقریبی تجربی} است، نه یک
\lr{PBE} گواهی‌شده. نظریه می‌گوید تعادل چه \emph{شکلی} باید داشته باشد
(ساختار آستانه‌ای، تفکیک از راه سیگنال)؛ خودبازی نزدیک آن را جست‌وجو می‌کند؛
و تورنمنت مقاومت آن را می‌سنجد.
\end{remark}

% ==================================================================== %
\section{گزاره‌های تحلیلی}\label{sec:props-stated}
% ==================================================================== %

چهار گزاره را از مدل استخراج می‌کنیم و برای هرکدام صریح می‌گوییم از کدام جزء
مدل آمده و با چه آماره‌ای قابل آزمون است.

\begin{keybox}
\textbf{گزاره‌ی ۱ (دروازه‌ی پایین).}
چون $B\gg X$ و هزینه‌ی جاری کوچک است، آستانه‌ی ادامه
$q^\ast=(g+\varphi-\Delta X)/(B-X)$ بسیار پایین است. پس خروج
\emph{داوطلبانه} تقریباً هیچ‌وقت بهینه نیست و خروج عقلانی، خروج \emph{اجباری}
است.\\[3pt]
\emph{از کجا آمده:} از $u_i$ (بخش ۹ متن) و بازه‌های $B$ و $X$ در شکل
رسمی‌شده.\quad
\emph{آماره:} توزیع $q^\ast$؛ نرخ خروج داوطلبانه؛ برتری استراتژی‌های
«هرگز خارج نشو».
\end{keybox}

\begin{keybox}
\textbf{گزاره‌ی ۲ (سیگنال بهینه درونی و غیریکنوا).}
سود نسبت به $\sig$ یکنوا نیست. $\sig=0.5$ بدترین انتخاب است چون هر مرحله
$\varphi$ می‌دهد بدون اینکه قفل نوبت را فعال کند؛ $\sig=0.75$ بهینه است؛
$\sig=1$ بدتر از $0.75$ است چون $\varphi\propto\sig^2$ و فرسایش $X$ بیشتر
است.\\[3pt]
\emph{از کجا آمده:} از قاعده‌ی نوبت‌دهی $M(h^t)$ + معافیت $\varphi$ در
\lr{HOLD} + شکل $\varphi$ و $X$.\quad
\emph{آماره:} توزیع $\sig$ انتخابی عامل؛ صورت‌حساب $\varphi$ به‌ازای هر سطح.
\end{keybox}

\begin{keybox}
\textbf{گزاره‌ی ۳ (خروج ساختنی است).}
چون $+20M^t$ بزرگ‌ترین تک‌جمله‌ی $X$ است، تصمیم‌های خروج باید حول پنجره‌ی
میانجی‌گری \emph{خوشه} بزنند.\\[3pt]
\emph{از کجا آمده:} از $X_i(t,h)$ (بخش ۷ متن).\quad
\emph{آماره:} نسبت $\Pr(\text{خروج}\mid M=1)$ به $\Pr(\text{خروج}\mid M=0)$.
\end{keybox}

\begin{keybox}
\textbf{گزاره‌ی ۴ (گیر افتادن در تعهد خود).}
چون $X$ در $K$ نزولی است، تمایل به خروج باید با انباشت تعهد عمومی خودِ
بازیگر \emph{کم} شود.\\[3pt]
\emph{از کجا آمده:} از جمله‌ی $-\lambda K$ در $X$.\quad
\emph{آماره:} $\Pr(\text{خروج})$ بر حسب سه‌بخشی‌های $K$.
\end{keybox}

% ==================================================================== %
\section{بخش شبیه‌سازی و پیاده‌سازی}\label{sec:impl}
% ==================================================================== %

\subsection{دو ایراد در کد ارائه‌شده}

\paragraph{(الف) ماژول گم‌شده — و اینکه چطور پیدایش کردیم.}
هم \lr{scripts/evaluate.py} و هم \lr{scripts/eval\_submission.py} از
\lr{see.tournament.runner} وارد می‌کنند، ولی این ماژول در درخت کاری مخزن
\emph{وجود ندارد} و هر دو اسکریپت با خطای \lr{ImportError} می‌افتند. ما اول
آن را طبق پروتکل مستندشده در \lr{docs/DESIGN.md} بخش ۱۰ بازسازی کردیم
(‏هر ورودی در برابر همه، \emph{در هر دو نقش}، روی فهرست بذر مشترک، امتیاز =
میانگین مطلوبیت خام با بازه‌ی اطمینان \lr{bootstrap}).

بعداً متوجه شدیم مخزن یک \emph{مخزن گیت} است و این فایل‌ها فقط از درخت کاری
\emph{حذف} شده‌اند، نه از تاریخچه. با
\lr{git checkout HEAD -- see/tournament/} نسخه‌ی اصلی برگشت — و همراهش
\lr{benchmarks/master\_agent.pt}، یعنی همان «بنچمارک بازی با استاد» که
دستورالعمل به آن اشاره می‌کند و بدون آن اصلاً قابل ارزیابی نبودیم.

مقایسه‌ی بازسازی ما با نسخه‌ی اصلی: پروتکل یکی است؛ تفاوت‌ها فقط در
جزئیات بذردهی است (‏\lr{seed0} پیش‌فرض $20{,}260{,}713$ به‌جای $7$، و
بذر عامل $s\cdot7919+\{1,2\}$ به‌جای $2s+\{0,1\}$). عددها هم عملاً یکی
درآمدند: امتیاز چک‌پوینت ارسالی $84.8$ با بازسازی ما و $84.2$ با ماژول
اصلی. \textbf{از اینجا به بعد همه‌ی اعداد گزارش با ماژول اصلی تولید
شده‌اند.} هیچ فیزیکی در \lr{config.py} یا \lr{env.py} دست نخورده است.

\begin{keybox}
\textbf{تایید بیرونی (۱۸ مرداد ۱۴۰۵).} استاد بعداً پوشه‌ی
\lr{tournament/} را در گروه کلاس منتشر کرد تا خطای
\lr{scripts/evaluate.py} رفع شود. آن را با نسخه‌ای که تمام مدت با آن کار
کرده‌ایم مقایسه کردیم (\lr{analysis/verify\_tournament.py}): پس از
یکسان‌سازی پایان‌خط، هر سه فایل \emph{بایت‌به‌بایت} یکی‌اند.

\vspace{3pt}
\centering\small
\begin{tabular}{lll}
\toprule
فایل & نسخه‌ی منتشرشده & نسخه‌ی ما\\
\midrule
\lr{\_\_init\_\_.py}   & \lr{92c10de525a9e226} & \lr{92c10de525a9e226}\\
\lr{runner.py}         & \lr{766b18a3a4deda31} & \lr{766b18a3a4deda31}\\
\lr{leaderboard.py}    & \lr{b9faec145a433d69} & \lr{b9faec145a433d69}\\
\bottomrule
\end{tabular}
\end{keybox}

نکته‌ی فنی: تفاوت خام فایل‌ها فقط \lr{CRLF} در برابر \lr{LF} است (نسخه‌ی
ما از \lr{checkout} گیت روی ویندوز آمده، فایل منتشرشده روی مک بسته‌بندی
شده). این تفاوت هر بایتِ هر خط را عوض می‌کند و هیچ رفتاری را عوض نمی‌کند،
پس مقایسه‌ی خام «کاملاً متفاوت» گزارش می‌دهد در حالی که فایل‌ها یکی‌اند.

فایل‌های منتشرشده را عمداً \emph{جایگزین نکردیم}. از نظر رفتاری یکی‌اند، و
خالی بودن خروجی \lr{git status --porcelain see/} همان شاهدی است که برای
بند ۵-۳ (دست‌نخوردگی موتور) ارائه می‌کنیم؛ رونویسی فایل‌ها این شاهد را
بی‌دلیل خراب می‌کرد.

\begin{answer}
\textbf{یک نکته‌ی صداقت.} تاریخچه‌ی گیت فایل‌های دیگری هم دارد که از درخت
کاری حذف شده‌اند، ولی همه‌شان \emph{سمت استاد}اند: وب‌اپلیکیشن
\lr{seeboard/}، پیکربندی‌های \lr{Dockerfile}/\lr{deploy/}، و
\lr{docs/handout/}. ما \emph{هیچ‌کدام} از این‌ها را برنگرداندیم و نگاه هم
نکردیم. فقط دو چیز را برگرداندیم که بدون‌شان ارزیابی ممکن نبود:
\lr{see/tournament/} (که اسکریپت‌های ارزیابیِ خودِ دانشجو بدون آن
\lr{ImportError} می‌دهند) و \lr{benchmarks/master\_agent.pt} (که
دستورالعمل صریحاً می‌گوید عامل در برابرش سنجیده می‌شود).
\end{answer}

\paragraph{(ب) تله‌ی ابرپارامتر — و خرابی‌ای که برمی‌گرداند.}

\begin{findbox}
\textbf{یافته‌ی کلیدی ۳.}
فایل \lr{docs/DESIGN.md} در بخش ۱۱-۳ یک خرابی به نام
«فروپاشی امتیازدهی خودبازی» را توصیف می‌کند (یک طرف یاد می‌گیرد در
$t\approx0$ خارج شود) و می‌گوید راه‌حلش «سهم حریف اسکریپتی $65\%$ و
آنتروپی $\ge0.02$» است. پیش‌فرض‌های خودِ \lr{PPOConfig} هم دقیقاً همین‌اند:
\lr{p\_self = 0.35} و \lr{entropy\_coef = 0.025}.

اما \lr{scripts/train.py} این‌ها را با پیش‌فرض‌های خط فرمان
\lr{--p-self 0.6} و \lr{--entropy 0.01} \textbf{بازنویسی می‌کند} — یعنی
اجرای دستور مستندشده‌ی پروژه، دقیقاً همان خرابی‌ای را فعال می‌کند که سند
می‌گوید رفع شده است.

اندازه‌گیری ما با \lr{analysis/probe.py} روی چک‌پوینت آموزش‌دیده با
پیش‌فرض‌های ارائه‌شده:
\begin{center}
\lr{as U.S.: exits in 100\% of episodes, mean t = 2.9,
first action $\sigma=0$ in 100\%}
\end{center}
یعنی صندلی آمریکا یک سیاست \emph{تسلیم محض} یاد گرفته بود.
\end{findbox}

راه‌حل ما: پاس دادن صریح مقادیر مستندشده
(\lr{--p-self 0.35 --entropy 0.025}). این یک پرچم \emph{مربی} است، نه تغییر
در \lr{see/}.

\subsection{تشخیص دوم: استخر حریفان بیش از حد فقیر است}

حتی بعد از اصلاح ابرپارامترها، صندلی آمریکا هنوز $100\%$ تسلیم می‌شد. علتش
ترکیب استخر حریف است: \lr{Random} خودش تصادفی خارج می‌شود و \lr{TitForTat}
وقتی $\ehat<0.20$ شود کوتاه می‌آید — یعنی \emph{هیچ حریفی در استخر نیست که
صرفاً حاضر نباشد برود}. پس یادگیرنده هیچ‌وقت برای زود تسلیم شدن تنبیه
نمی‌شود.

ما پنج الگو در \lr{my\_team/opponents.py} ساختیم
(‏\lr{Dove}، \lr{Hawk}، \lr{Plateau}، \lr{Patient}، \lr{BayesThreshold}) و
با \lr{my\_team/train\_plus.py} آن‌ها را \emph{در حافظه} به استخر مربی تزریق
کردیم. هیچ فایلی زیر \lr{see/} تغییر نکرده؛ فقط \emph{با چه کسی تمرین
می‌کنیم} عوض شده — و \lr{DESIGN.md} بخش ۱۱ صریحاً می‌گوید همین پیچ برای
تنظیم گذاشته شده است. اثر فوری: طول اپیزود از $3.9$ به $9.2$ پرید.

\begin{answer}
\textbf{دقیقاً چه چیزی عوض شده و چه چیزی نه.} بند ۵-۳ دستورالعمل می‌گوید
معماری شبکه و الگوریتم آموزش برای همه یکسان است و نباید تغییر کند. پس صریح
می‌نویسم:
\begin{itemize}[leftmargin=1.3em,itemsep=1pt]
  \item \textbf{دست نخورده:} کل پوشه‌ی \lr{see/} (فیزیک بازی، محیط،
        پیکربندی)، معماری شبکه (\lr{GRU} + سرهای سیاست و ارزش)، و خود
        الگوریتم \lr{PPO}. یک بایت هم عوض نشده است.
  \item \textbf{عوض شده، و همه‌شان پرچم مربی‌اند:} مقدار \lr{p\_self}،
        مقدار \lr{entropy\_coef} (که در طول آموزش از $0.02$ به $0.004$
        کم می‌شود)، تعداد گام‌ها، و اینکه چه حریفانی در استخر تمرین باشند.
        هر چهارتا در \lr{scripts/train.py} هم به‌صورت پرچم خط فرمان وجود
        دارند؛ ما فقط مقدارشان را عوض کرده‌ایم.
  \item \textbf{در چک‌پوینت نهایی:} استخر تمرین \lr{--pool stock} است،
        یعنی \emph{فقط} همان دو حریف مرجع. پنج الگوی بالا در مسیر تشخیص
        به‌کار آمدند ولی در آموزش نسخه‌ی ارسالی استفاده نشده‌اند.
\end{itemize}
پس تفاوت عملکرد از نظریه می‌آید (ویژگی‌ها و شکل‌دهی در \lr{theory.py})، نه
از مهندسی — همان چیزی که دستورالعمل می‌خواهد.
\end{answer}

\subsection{نظریه‌ی ما: فیلتر بیزی روی نوع حریف}

فایل \lr{my\_team/theory.py} دو تابع دارد.

\paragraph{تابع اول — ویژگی‌های باور.}
بر پایه‌ی مشاهده‌ی بخش~\ref{sec:beliefs}، یک شبکه‌ی وزن‌دار
$11\times9=99$ نقطه‌ای روی $\theta_j=(\rbar_j,\mbar_j)$ نگه می‌داریم. هر
مرحله: (۱) تاب‌آوری هر ذره را با قانون \emph{عمومی} هزینه جلو می‌بریم؛
(۲) ذره‌هایی که تاب‌آوری‌شان به صفر رسیده ولی حریف هنوز نرفته را
\emph{حذف} می‌کنیم (کانال بقا)؛ (۳) با
$\exp(-\beta\varphi_j)$ وزن‌دهی می‌کنیم (کانال تک‌تقاطعی). از این پسین
۱۴ ویژگی می‌سازیم، از جمله $\hat q$ (احتمال بردن مسابقه، با استفاده از فرمول
بسته‌ی گزاره~\ref{pr:collapse}) و خودِ دروازه‌ی گزاره~\ref{pr:gate}.
بعداً یک کانال چهارم و دو ویژگی دیگر اضافه شد (مجموعاً ۱۶)؛ چرایی‌اش در
بخش~\ref{sec:tourn} توضیح داده شده، چون از دل یک اندازه‌گیری بیرون آمد نه از
دل متن.

\paragraph{اعتبارسنجی فیلتر.} با \lr{analysis/filter\_check.py} پسین را با
مقدار واقعی (که عامل هرگز نمی‌بیند) مقایسه کردیم:

\begin{table}[h]\centering
\begin{tabular}{lcccc}
\toprule
کمیت & بایاس & \lr{MAE} & \lr{RMSE} & \lr{MAE} پیش‌بینی ثابت\\
\midrule
$\ehat_j$ (تاب‌آوری پنهان حریف) & $+0.0033$ & $\mathbf{0.0511}$ & $0.0835$ & $0.1884$\\
$\rbar_j$ (عزم پایه‌ی حریف)     & $-0.0154$ & $0.1522$ & $0.1836$ & ---\\
\bottomrule
\end{tabular}
\caption{فیلتر عملاً بدون بایاس است و خطای برآورد تاب‌آوری پنهان را
\textbf{$\mathbf{72.9\%}$} نسبت به پیش‌بینی ثابت کم می‌کند
(‏$7438$ مشاهده، $250$ اپیزود).}
\end{table}

\paragraph{تابع دوم — شکل‌دهی پاداش.}
تابع پتانسیل را برابر \emph{ارزش دو گزینه‌ی پایانی‌ای که همین حالا در دست
داریم} گذاشتیم:
\[
  \Phi(s)=0.5\,X_i(t,h^t)+0.10\,B_i(e_i,\rbar_i),
  \qquad F=\gamma\Phi(s')-\Phi(s).
\]
منطقش: هر دو جمله فقط در \emph{آخرین} مرحله پرداخت می‌شوند، پس \lr{PPO}
تعارض اعتبار-انعطاف‌پذیری را فقط از یک عدد خیلی تأخیری می‌بیند. تفاضل‌گیری،
\emph{فرسایش} گزینه‌ی خروج را همان مرحله‌ای که انتخاب انجام می‌شود
صورت‌حساب می‌کند. چون شکل‌دهی \emph{پتانسیل‌پایه} است، در طول اپیزود جمله‌هایش
همدیگر را حذف می‌کنند و
سیاست بهینه را تغییر نمی‌دهد (‏\lr{Ng, Harada \& Russell 1999}) — پس فقط
می‌تواند یادگیری را تند کند، نه هدف را منحرف.

\subsection{فضای استراتژی‌ها: یک جست‌وجوی منظم}\label{sec:sweep}

قبل از اینکه به عامل آموزش‌دیده اعتماد کنیم، فضای استراتژی‌های ساده را
یکی‌یکی بررسی کردیم (\lr{analysis/sweep.py})، در همان تنظیمی که
\lr{scripts/evaluate.py} استفاده می‌کند (شما در برابر \lr{Random} و
\lr{TitForTat}). هر استراتژی یک زوج (سطح سیگنال ثابت، آستانه‌ی خروج) است:

\begin{table}[h]\centering
\begin{tabular}{lrrr}
\toprule
استراتژی & امتیاز خودش & امتیاز \lr{TitForTat} & اختلاف\\
\midrule
$\sig=0.75$، هرگز خارج نشو            & $47.6$ & $30.2$ & $\mathbf{+17.4}$\\
$\sig=1.0$، هرگز خارج نشو             & $34.4$ & $20.6$ & $+13.8$\\
آینه‌ی حریف با سقف $0.75$، هرگز خارج نشو & $\mathbf{67.9}$ & $56.3$ & $+11.6$\\
\midrule
$\sig=0.75$، خروج در $\ehat<0.35$      & $36.1$ & $59.8$ & $-23.7$\\
$\sig=0.0$، هرگز خارج نشو             & $11.7$ & $56.3$ & $-44.6$\\
$\sig=0.5$، هرگز خارج نشو             & $-10.6$ & $39.9$ & $-50.4$\\
\bottomrule
\end{tabular}
\caption{گزیده‌ای از ۲۶ استراتژی بررسی‌شده. دو الگو کاملاً روشن است.}
\end{table}

\begin{answer}
\textbf{هر دو گزاره‌ی ۱ و ۲ اینجا تایید می‌شوند.}
\begin{itemize}[leftmargin=1.3em,itemsep=1pt]
  \item \textbf{گزاره‌ی ۱:} هر استراتژی با «هرگز خارج نشو» بین $20$ تا $50$
        واحد بهتر از دوقلوی خودش با آستانه‌ی خروج است.
  \item \textbf{گزاره‌ی ۲:} ترتیب $\sig$ عبارت است از
        $0.75>1.0>0.0>0.5$ — \emph{کاملاً غیریکنوا}، دقیقاً همان‌طور که
        استدلال معافیت $\varphi$ در \lr{HOLD} پیش‌بینی می‌کرد. $\sig=0.5$
        بدترین است.
\end{itemize}
\end{answer}

\subsection{اتصال نظریه به رفتار مشاهده‌شده}\label{sec:props}

حالا چهار گزاره را روی داده‌ی تولیدشده توسط \emph{عامل آموزش‌دیده}
می‌سنجیم (\lr{analysis/propositions.py}).

همه‌ی اعداد این بخش از خودِ چک‌پوینت ارسالی (\lr{submission.pt}) روی
$400$ اپیزود در هر جفت و $7{,}730$ مرحله‌ی تصمیم به دست آمده‌اند.

\begin{answer}
\textbf{یک اشکال در آماره‌ی خودمان که وسط کار پیدا شد.} اولین نسخه‌ی
\lr{propositions.py} هر مرحله‌ای را که در آن \lr{EXIT} مجاز بود «فرصت خروج
داوطلبانه» می‌شمرد. ولی وقتی تاب‌آوری تمام می‌شود موتور \emph{فقط}
\lr{EXIT} را مجاز می‌گذارد؛ آن خروج اجباری است، نه انتخاب. این خطا همه‌ی
آماره‌های زیر را جابه‌جا می‌کرد و مخصوصاً \emph{علامت} گزاره‌ی ۴ را برعکس
نشان می‌داد، چون خروج اجباری دقیقاً در مراحل دیر رخ می‌دهد، یعنی همان جایی
که $K$ خودِ بازیکن بیشترین مقدار را دارد. با شرط «ماندن هم مجاز بود» آماره
درست شد. عددهای زیر همه با شرط اصلاح‌شده‌اند.
\end{answer}

\paragraph{گزاره‌ی ۱ — دروازه‌ی ادامه.}
میانه‌ی $q^\ast$ تجربی برابر $\mathbf{0.042}$ است — یعنی عامل باید ادامه دهد
مگر آنکه شانس بردنش زیر $4\%$ باشد. و واقعاً همین کار را می‌کند: از
$7{,}730$ مرحله‌ای که خروج \emph{و ماندن هر دو} مجاز بوده، فقط در
$\mathbf{6.9\%}$ خارج شده است. رفتار مشروط به نوع هم دقیقاً همان ساختار
آستانه‌ای گزاره~\ref{pr:cutoff} را نشان می‌دهد:

\begin{table}[h]\centering
\begin{tabular}{lcccc}
\toprule
چارک عزم پایه $\rbar$ & \lr{Q1} (کم) & \lr{Q2} & \lr{Q3} & \lr{Q4} (زیاد)\\
\midrule
$\Pr(\text{خروج})$ & $0.583$ & $0.443$ & $0.343$ & $\mathbf{0.165}$\\
میانگین زمان خروج & $4.1$ & $4.4$ & $4.7$ & $\mathbf{6.3}$\\
مطلوبیت میانگین $u$ & $46.6$ & $62.2$ & $69.9$ & $\mathbf{86.7}$\\
\bottomrule
\end{tabular}
\caption{احتمال خروج \emph{یکنواخت نزولی} در عزم پایه، و زمان خروج و
مطلوبیت هر دو \emph{یکنواخت صعودی} — یعنی عامل واقعاً رفتار مشروط به نوع
یاد گرفته است، با آنکه نوع خودش را فقط از دو عدد $(\rbar,\mhat)$ در بردار
مشاهده می‌بیند.}
\end{table}

\paragraph{گزاره‌ی ۲ — سیگنال بهینه درونی.}

\begin{table}[h]\centering
\begin{tabular}{lccccc}
\toprule
$\sig$ & $0.00$ & $0.25$ & $0.50$ & $0.75$ & $1.00$\\
\midrule
سهم از مراحل & $\mathbf{24.8\%}$ & $2.2\%$ & $0.3\%$ & $\mathbf{72.4\%}$ & $0.3\%$\\
$\varphi$ میانگین & $0.000$ & $0.137$ & $0.620$ & $1.312$ & $2.992$\\
\bottomrule
\end{tabular}
\caption{$97.2\%$ از مراحل روی \emph{فقط دو} سطح $\{0,\,0.75\}$ جمع شده و
سه سطح دیگر عملاً خالی‌اند. نسبت این دو سطح نسبت به نسخه‌های قبلی برعکس
شده (‏$72.4\%$ روی $0.75$ به‌جای $17\%$)، چون استخر آموزش تازه پر از
حریفی است که کوتاه نمی‌آید.}
\end{table}

\begin{answer}
توزیع سیگنال \emph{یکنوا نیست} و شکلش دقیقاً همان چیزی است که تحلیل
پیش‌بینی کرده بود: یا زیر آستانه بمان و اصلاً پول سیگنال نده، یا تا خود
آستانه‌ی قفل برو. سطح میانی $0.5$ بدترین گزینه است چون هر مرحله $\varphi$
می‌دهد بی‌آنکه قفل نوبت را فعال کند، و $\sig=1$ فقط $\varphi$ را گران‌تر
می‌کند بدون مزیت تازه. این همان تعادل \emph{نیمه‌تجمیعی} پیش‌بینی‌شده در
بخش تحلیل تعادلی است، و در این چک‌پوینت \emph{صریح‌تر} از قبل ظاهر شده:
توده‌ی احتمال دقیقاً روی دو سطحی نشسته که تحلیل نام برده بود.

یک نکته‌ی اضافی که خودش از داده درآمد: شدت سیگنال با \emph{کم شدن}
تاب‌آوری بالا می‌رود ($\Pr(\sig\ge0.75)$ برابر $0.944$ در سه‌بخشی پایینِ
تاب‌آوری در برابر $0.782$ در سه‌بخشی بالا). یعنی عامل موضع بلند را برای وقتی نگه
می‌دارد که چیزی برای از دست دادن ندارد — رفتاری که در مدل جایی برایش
پیش‌بینی نکرده بودیم.
\end{answer}

\paragraph{گزاره‌ی ۳ — خروج ساختنی.}
\begin{table}[h]\centering
\begin{tabular}{lcc}
\toprule
 & پنجره‌ی میانجی‌گری باز & بسته\\
\midrule
$\Pr(\text{خروج داوطلبانه}\mid\text{ماندن هم مجاز})$ & $\mathbf{0.0802}$ & $0.0670$\\
تعداد مرحله & $910$ & $6{,}820$\\
\bottomrule
\end{tabular}
\caption{نسبت $\mathbf{1.20}$ برابر — جهت درست است ولی اثر بسیار ضعیف‌تر
از نسخه‌های قبلی که $3.4$ برابر می‌داد.}
\end{table}

\begin{answer}
\textbf{چرا این گزاره در چک‌پوینت نهایی ضعیف شده — و چرا خودِ مدل جوابش را
دارد.} پنجره‌ی میانجی‌گری فقط وقتی با احتمال $0.35$ باز می‌شود که موضع
مشترک آرام باشد ($\tilde\sig_I+\tilde\sig_U\le0.5$)؛ وگرنه احتمالش
$0.05$ است. چک‌پوینت نهایی $72.4\%$ مراحل را روی $\sig=0.75$ می‌ماند، پس
عملاً \emph{خودش} پنجره را بسته نگه می‌دارد: از $7{,}730$ مرحله فقط
$910$ تا با پنجره‌ی باز اتفاق می‌افتد، در حالی که در نسخه‌ی آرام‌تر قبلی
$12{,}198$ از $27{,}170$ بود.

پس این یک نقض گزاره‌ی ۳ نیست، بلکه \emph{حلقه‌ی میانجی} جمع‌بندی است که
روی داده دیده می‌شود: تشدید $\Rightarrow$ پنجره بسته می‌ماند $\Rightarrow$
خروج آبرومندانه در دسترس نیست. عامل این معامله را آگاهانه انجام می‌دهد،
چون در استخر تازه بیشترِ حریفان کسانی‌اند که با آرامش نمی‌شود از پا
درشان آورد.
\end{answer}

\paragraph{گزاره‌ی ۴ — گیر افتادن در تعهد خود.}
\begin{table}[h]\centering
\begin{tabular}{lccc}
\toprule
انبار تعهد $K$ & کم ($0.44$) & متوسط ($1.33$) & زیاد ($2.28$)\\
\midrule
$\Pr(\text{خروج داوطلبانه})$ & $0.1371$ & $0.0319$ & $\mathbf{0.0036}$\\
تعداد مرحله & $3{,}370$ & $1{,}850$ & $2{,}510$\\
\bottomrule
\end{tabular}
\caption{تمایل به خروج با انباشت تعهد عمومی \emph{یکنواخت} و در مجموع حدود
\emph{سی‌وهشت برابر} کم می‌شود — دقیقاً پیش‌بینی جمله‌ی $-\lambda K$. و چون
این عامل $K$ را عمداً بالا می‌برد، بازه‌ی $K$ این بار پهن است و شیب واقعاً
دیده می‌شود.}
\end{table}

با \lr{analysis/k\_control.py} همین جدول را داخل بازه‌های ثابت مرحله و
تاب‌آوری باقی‌مانده هم حساب کردیم، چون $K$ خودش با گذشت زمان بالا می‌رود و
ممکن بود اثر واقعی از آنِ زمان باشد نه $K$:

\begin{table}[h]\centering\small
\begin{tabular}{lcccc}
\toprule
 & $t=0..4$ & $t=5..9$ & $t=10..14$ & $t\ge15$\\
\midrule
$\Pr(\text{خروج})$، $K$ زیر میانه   & $0.1170$ & $0.0105$ & --- & ---\\
$\Pr(\text{خروج})$، $K$ بالای میانه & $\mathbf{0.0272}$ & $\mathbf{0.0000}$ & $\mathbf{0.0000}$ & $\mathbf{0.0000}$\\
\bottomrule
\end{tabular}
\caption{برش روی \emph{میانه}‌ی $K$ همین سیاست گرفته شده ($1.31$)، نه یک
عدد ثابت — با یک برش ثابت تقریباً همه‌ی مراحل در یک سبد می‌افتادند. در هر
بازه‌ی زمانی (و در آزمون دوم، در هر بازه‌ی تاب‌آوری باقی‌مانده) نرخ خروج
با $K$ بالاتر کمتر است. پس اثر واقعاً از $K$ می‌آید، نه از همبستگی $K$ با
دیر بودن مرحله.}
\end{table}

\begin{answer}
\textbf{هر چهار گزاره تایید می‌شوند}، و هر چهار آماره در جهتی است که تحلیل
بخش‌های ۴ تا ۸ پیش‌بینی کرده بود.

یک تذکر لازم درباره‌ی اینکه این تاییدها چقدر «مستقل» از خودِ ما هستند.
جمله‌ی اصلی شکل‌دهی \emph{پتانسیل‌پایه} است و طبق قضیه‌ی
\lr{Ng, Harada \& Russell} سیاست بهینه را عوض نمی‌کند، و ویژگی‌های باور هم
فقط \emph{آماره} می‌دهند نه دستور. پس گزاره‌های ۲ تا ۴ — که همه درباره‌ی
\emph{ساختار شرطی} رفتارند (سیگنال بر حسب سطح، خروج بر حسب پنجره، خروج بر
حسب $K$) — چیزی‌اند که عامل خودش پیدا کرده است.

و در چک‌پوینت \emph{نهایی} این تذکر حتی لازم هم نیست: جمله‌ی دوم شکل‌دهی
— تنها چیزی که در این نظریه پتانسیل‌پایه نبود — روی صفر تنظیم شده است
(بخش~\ref{sec:realboard}). پس کل شکل‌دهی پتانسیل‌پایه است و طبق آن قضیه
سیاست بهینه را تغییر نمی‌دهد. یعنی \emph{هیچ‌کدام} از چهار آماره‌ی بالا را
ما به عامل تحمیل نکرده‌ایم؛ همه از دل بیشینه‌کردن همان مطلوبیت خام مدل
بیرون آمده‌اند.
\end{answer}

\subsection{تشخیص سوم: استخر آموزش باید با معیار ارزیابی هماهنگ باشد}

استخر غنی‌شده طول اپیزود را درست کرد ولی امتیاز را نه. وقتی دلیلش را
بررسی کردیم، به نکته‌ای رسیدیم که در نگاه اول خلاف انتظار است.

\begin{findbox}
\textbf{یافته‌ی کلیدی ۴ — چرا «تسلیم زودهنگام» گاهی درست است.}
فرض کنید حریف \emph{هرگز} داوطلبانه خارج نمی‌شود (مثل \lr{Dove}) یا آستانه‌ی
خروجش آن‌قدر پایین است که در افق $\bar T=40$ به آن نمی‌رسد (مثل
\lr{TitForTat} با $\ehat<0.20$ در موضع آرام). آنگاه بازی به \emph{مهلت}
می‌رسد و هر دو $X_i(\bar T)$ می‌گیرند — ولی $G_i$ در آن نقطه تقریباً برابر
کل $e_i^0$ است:
\[
  u_i^{\text{مهلت}}\approx X_i-G_i(\bar T)\approx 28-53<0,
  \qquad\text{در برابر}\qquad
  u_i^{\text{خروج در }t=3}\approx X_i-G_i(3)\approx 28-4 .
\]
پس در برابر حریفی که نمی‌رود، \textbf{خروج زودهنگام واقعاً بهینه است} — آنچه
ما اول «خرابی تسلیم» تشخیص داده بودیم، در برابر بخشی از استخر یک
\emph{بهترین پاسخ} بود.

و این توضیح می‌دهد چرا استخر غنی‌شده کمکی نکرد: ما با اضافه‌کردن
\lr{Dove}/\lr{Patient}/\lr{Hawk} تعداد حریفانی را که \emph{نمی‌روند}
\emph{زیاد} کردیم، پس خروج زودهنگام را \emph{جذاب‌تر} کردیم.
\end{findbox}

راه‌حل: استخر آموزش را با همان توزیعی هماهنگ کنیم که در آن ارزیابی می‌شویم
(‏\lr{Random} و \lr{TitForTat})، با
\lr{--pool stock --p-self 0.10 --p-scripted 1.0}. اثرش چشمگیر بود:

\begin{table}[h]\centering
\begin{tabular}{lccc}
\toprule
رژیم آموزش & طول اپیزود & $u_I$ & $u_U$\\
\midrule
پیش‌فرض‌های ارائه‌شده & $3.2$ & $144.8$ & $30.8$\\
ابرپارامترهای مستندشده & $3.9$ & $129.7$ & $33.4$\\
استخر غنی‌شده & $5.7$ & $110.9$ & $35.1$\\
\textbf{هماهنگ با مرجع} & $\mathbf{10.4}$ & $\mathbf{89.6}$ & $\mathbf{70.0}$\\
\bottomrule
\end{tabular}
\caption{آمار پایان آموزش. در رژیم آخر دو نقش تقریباً \emph{متوازن}
می‌شوند — یعنی صندلی آمریکا دیگر یک سیاست تسلیم محض نیست.}
\end{table}

\subsection{نتایج ارزیابی و مطالعه‌ی حذفی}

\begin{table}[h]\centering
\begin{tabular}{lrrr}
\toprule
پیکربندی & امتیاز کل & به‌عنوان ایران & به‌عنوان آمریکا\\
\midrule
\lr{TitForTat} (مرجع) & $95.5$ & $152.7$ & $38.2$\\
\midrule
\textbf{ویژگی + شکل‌دهی، هماهنگ با مرجع} (ارسالی)
                          & $\mathbf{89.8}$ & $136.3$ & $43.4$\\
فقط شکل‌دهی، هماهنگ با مرجع & $89.8$ & $135.5$ & $44.1$\\
ویژگی + شکل‌دهی، استخر غنی‌شده & $84.5$ & $124.8$ & $44.2$\\
فقط شکل‌دهی، استخر غنی‌شده & $75.8$ & $106.9$ & $44.7$\\
فقط شکل‌دهی، پیش‌فرض‌های ارائه‌شده & $69.8$ & $91.7$ & $47.9$\\
\lr{Random} (مرجع) & $38.4$ & $64.2$ & $12.5$\\
\bottomrule
\end{tabular}
\caption{تورنمنت دوره‌ای در میدان گسترده ($30$ اپیزود در هر جفت مرتب،
بازه‌ی اطمینان $95\%$ حدود $\pm8$). بازه‌ی چک‌پوینت ارسالی
$[82.0,97.2]$ با بازه‌ی \lr{TitForTat} یعنی $[86.5,104.3]$
\emph{هم‌پوشانی} دارد.}
\end{table}

\subsection{ارزیابی در برابر بنچمارک استاد}

بازیابی \lr{benchmarks/master\_agent.pt} از تاریخچه‌ی گیت اجازه داد آزمون
واقعی را انجام دهیم. نتیجه با ماژول تورنمنت \emph{اصلی} و $60$ اپیزود در
هر جفت مرتب:

\begin{table}[h]\centering
\begin{tabular}{lrrrl}
\toprule
عامل & امتیاز & ایران & آمریکا & بازه‌ی اطمینان $95\%$\\
\midrule
\lr{TitForTat} (مرجع) & $81.3$ & $115.0$ & $47.6$ & $[71.2,\ 90.2]$\\
\textbf{نسخه‌ی ۱} & $\mathbf{72.4}$ & $84.8$ & $60.0$ & $[64.4,\ 80.2]$\\
\lr{MasterAgent} (بنچمارک استاد) & $38.5$ & $40.3$ & $36.6$ & $[30.6,\ 45.7]$\\
\lr{Random} (مرجع) & $9.3$ & $9.5$ & $9.2$ & $[3.9,\ 14.7]$\\
\bottomrule
\end{tabular}
\end{table}

\begin{findbox}
\textbf{چک‌پوینت ارسالی بنچمارک استاد را با اختلاف زیاد می‌برد}
(‏$72.4$ در برابر $38.5$؛ بازه‌های اطمینان \emph{هم‌پوشانی ندارند})، و در
رویارویی مستقیم \emph{در هر دو نقش} جلو است:
به‌عنوان ایران $70.4$ در برابر $17.1$، و به‌عنوان آمریکا $23.4$ در برابر
$6.2$.
\end{findbox}

\paragraph{تنها نقطه‌ی ضعف، و اینکه چرا هست.}
رویارویی‌های مستقیم چک‌پوینت ارسالی:

\begin{table}[h]\centering\small
\begin{tabular}{llrrr}
\toprule
ایران & آمریکا & $u_I$ & $u_U$ & طول\\
\midrule
\textbf{نسخه‌ی ۱} & \lr{Random} & $\mathbf{127.6}$ & $-4.6$ & $6.1$\\
\textbf{نسخه‌ی ۱} & \lr{MasterAgent} & $\mathbf{70.4}$ & $17.1$ & $17.7$\\
\textbf{نسخه‌ی ۱} & \lr{TitForTat} & $\mathbf{56.5}$ & $-3.3$ & $24.7$\\
\lr{Random} & \textbf{نسخه‌ی ۱} & $2.3$ & $\mathbf{113.1}$ & $6.2$\\
\lr{MasterAgent} & \textbf{نسخه‌ی ۱} & $6.2$ & $\mathbf{23.4}$ & $10.4$\\
\lr{TitForTat} & \textbf{نسخه‌ی ۱} & $159.0$ & $43.7$ & $4.6$\\
\bottomrule
\end{tabular}
\caption{در پنج رویارویی از شش، نسخه‌ی ۱ حریف را می‌برد — از جمله
\emph{هر دو} رویارویی با بنچمارک استاد و رویارویی با \lr{TitForTat} وقتی
خودمان ایران هستیم ($56.5$ در برابر $-3.3$).}
\end{table}

تنها شکست، سطر آخر است: وقتی ما آمریکا هستیم و \lr{TitForTat} ایران، ما در
$t\approx4.6$ خارج می‌شویم و $43.7$ می‌گیریم، در حالی که \lr{TitForTat}
جایزه‌ی دوام‌آوردن $159.0$ را برمی‌دارد. و این دقیقاً همان محاسبه‌ی
یافته‌ی کلیدی ۴ است: اگر مقاومت می‌کردیم، بازی تا $\sim25$ مرحله می‌رفت و
سهم \emph{خودمان} از $43.7$ به حدود $15$ می‌افتاد (چون $G$ انباشته
می‌شود)، هرچند سهم \lr{TitForTat} هم از $159$ به نزدیک صفر می‌رسید.

\begin{answer}
\textbf{گزارش صادقانه‌ی جایگاه — در این مرحله از کار.} نسخه‌ی ۱ بنچمارک
استاد را قاطعانه می‌برد ولی در امتیاز کل هنوز پشت \lr{TitForTat} است
(‏$72.4$ در برابر $81.3$، با بازه‌های اطمینان \emph{هم‌پوشان}). کل این
فاصله از یک رویارویی می‌آید. اینجا یک تنش واقعی بین دو هدف هست: بیشینه‌کردن
\emph{مطلوبیت خودمان} (که خروج زودهنگام را می‌طلبد) و
\emph{محروم‌کردن حریف} از جایزه (که مقاومت را می‌طلبد). چون امتیاز تورنمنت
کلاسی میانگین مطلوبیت \emph{خودِ} عامل روی همه‌ی حریفان است و
\lr{TitForTat} فقط یکی از آن‌هاست، ما هدف اول را انتخاب کردیم. اگر معیار
دقیقاً «شکست دادن \lr{TitForTat} در یک استخر سه‌تایی» بود، انتخاب بهینه
عوض می‌شد — و جست‌وجوی بخش~\ref{sec:sweep} نشان می‌دهد استراتژی‌های ساده‌ای
هست که با اختلاف $+11$ تا $+17$ از \lr{TitForTat} جلو می‌زنند، دقیقاً چون
حاضرند مطلوبیت خودشان را فدا کنند.

بخش بعدی همین را دنبال می‌کند و نشان می‌دهد معیار نمره‌دهی \emph{دقیقاً}
همان معیار دوم است، نه اولی. نتیجه‌ی نهایی این تعقیب در
بخش~\ref{sec:both} آمده: چک‌پوینت ارسالی در همان ارزیابی رسمی $71.8$ در
برابر $61.1$ می‌گیرد.
\end{answer}

\subsection{معیار تورنمنت: چرا بیشینه‌کردن مطلوبیت خود کافی نیست}
\label{sec:tourn}

\begin{answer}
\textbf{هشدار به خواننده.} تمام این زیربخش و سه زیربخش بعدی روی استخر
\emph{سه‌تایی} \lr{scripts/evaluate.py} بنا شده‌اند. آن استخر معیار
نمره‌دهی نیست. تحلیل زیر برای آن استخر \emph{درست} است و ما نگهش داشته‌ایم
چون مسیر فکری واقعی همین بود، ولی نتیجه‌گیری نهایی‌اش را
بخش~\ref{sec:realboard} پس می‌گیرد. اگر فقط دنبال نظریه‌ی ارسالی هستید،
مستقیم به آنجا بروید.
\end{answer}

جدول بالا یک چیز را روشن می‌کند: بنچمارک استاد را می‌بریم ولی از
\lr{TitForTat} عقبیم. سوال بعدی این بود که \emph{دقیقاً کجا} عقبیم. برای
جواب دادن، معیار را از اول باز کردیم.

\paragraph{معیار به دو مسئله‌ی مستقل تجزیه می‌شود.}
در ارزیابی رسمی استخر $\{\text{ما},\ \mathrm{Random},\ \mathrm{TitForTat}\}$ است و
هیچ عاملی با \emph{خودش} بازی نمی‌کند. پس شبکه‌ی ایران ما و شبکه‌ی آمریکای ما
دو مجموعه بازی \emph{مجزا} می‌بینند. با نمادگذاری
\[
\begin{aligned}
A&=u_I(\text{ما ایران}\mid \mathrm{Rnd}), &\quad
B&=u_I(\text{ما ایران}\mid \mathrm{TFT}), &\quad
b'&=u_U(\mathrm{TFT}\mid \text{ما}),\\
C&=u_U(\text{ما آمریکا}\mid \mathrm{Rnd}), &\quad
D&=u_U(\text{ما آمریکا}\mid \mathrm{TFT}), &\quad
d'&=u_I(\mathrm{TFT}\mid \text{ما}),
\end{aligned}
\]
و $R_1=u_I(\mathrm{TFT}\mid\mathrm{Rnd})$، $R_2=u_U(\mathrm{TFT}\mid\mathrm{Rnd})$،
دو امتیاز جدول برابرند با $(A{+}B{+}C{+}D)/4$ و
$(R_1{+}R_2{+}b'{+}d')/4$. پس معیار «امتیاز بالاتر از \lr{TitForTat}»
\emph{دقیقاً} این است:
\begin{equation}\label{eq:crit}
  \underbrace{(A+B-b')}_{J_I\ :\ \text{فقط صندلی ایران}}
  \;+\;
  \underbrace{(C+D-d')}_{J_U\ :\ \text{فقط صندلی آمریکا}}
  \;>\; R_1+R_2 .
\end{equation}
$R_1$ و $R_2$ اصلاً به ما بستگی ندارند. اندازه‌گیری
(\lr{analysis/frontier.py}، ۶۰ اپیزود): $R_1=127.8$، $R_2=104.1$، پس آستانه
$231.9$ است. جست‌وجوی ۴۷ موضع مختلف روی هر صندلی:

\begin{table}[h]\centering\small
\begin{tabular}{lrrrr}
\toprule
موضع صندلی ایران & $A$ & $B$ & $b'$ & $J_I$\\
\midrule
\textbf{چک‌پوینت (ما)} & $129.1$ & $56.7$ & $-8.5$ & $\mathbf{194.3}$\\
$\sig=0.75$، هرگز خارج نشو & $133.9$ & $-0.5$ & $-58.8$ & $192.2$\\
$\sig=0$ با سقف $0.75$، هرگز خارج نشو & $131.6$ & $35.4$ & $-10.1$ & $177.0$\\
\midrule
موضع صندلی آمریکا & $C$ & $D$ & $d'$ & $J_U$\\
\midrule
\textbf{چک‌پوینت (ما)} & $108.1$ & $\mathbf{42.9}$ & $\mathbf{156.8}$ & $\mathbf{-5.8}$\\
$\sig=1.0$، هرگز خارج نشو & $97.3$ & $-50.6$ & $-65.1$ & $111.8$\\
$\sig=0.75$، هرگز خارج نشو & $109.2$ & $-37.3$ & $-38.1$ & $110.0$\\
\bottomrule
\end{tabular}
\caption{صندلی ایران ما \emph{بهترین} چیزی است که پیدا کردیم
($J_I=194.3$، بالاتر از هر موضع دستی). کل کسری در $J_U$ است: چک‌پوینت
بیشترین $D$ را می‌گیرد ($42.9$) و دقیقاً به همین دلیل بیشترین $d'$ را هم
\emph{هدیه می‌دهد} ($156.8$).}
\end{table}

\begin{findbox}
\textbf{یافته‌ی کلیدی ۵.}
عامل ما در صندلی آمریکا کار \emph{اشتباهی} نمی‌کند: $D=42.9$ در برابر
سقف نظری $X_U^{\max}=48$، یعنی تقریباً بیشترین چیزی که آن صندلی می‌تواند
بگیرد. مشکل این است که همان اقدامِ بهینه، جایزه‌ی $B_j\approx160$ را کامل
تحویل حریف می‌دهد. \textbf{بیشینه‌کردن $u_i$ و بیشینه‌کردن $u_i-u_j$ اینجا
در دو جهت مخالف‌اند.}
\end{findbox}

\paragraph{و این اثر استخر کوچک نیست.} اولین حدس ما این بود که چون استخر
مرجع فقط سه‌تایی است، یک رویارویی وزن غیرعادی دارد و در تورنمنت کلاسی
با $N$ شرکت‌کننده کم‌رنگ می‌شود. با \lr{analysis/classpool.py} این را آزمودیم
(چک‌پوینت‌های دیگرمان و بنچمارک استاد به‌جای هم‌کلاسی‌ها):

\begin{table}[h]\centering\small
\begin{tabular}{rrrrr}
\toprule
اندازه‌ی استخر & امتیاز ما & (ایران / آمریکا) & \lr{TitForTat} & اختلاف\\
\midrule
$3$ & $86.5$ & $96.8$ / $76.1$ & $96.9$ & $-10.5$\\
$5$ & $82.3$ & $104.8$ / $59.8$ & $95.2$ & $-12.9$\\
$8$ & $89.8$ & $128.5$ / $51.1$ & $99.4$ & $-9.6$\\
\bottomrule
\end{tabular}
\caption{اختلاف عملاً ثابت می‌ماند. دلیلش این است که \emph{هر} صندلی
آمریکا در استخر همین کار را می‌کند، پس \lr{TitForTat} از همه باج می‌گیرد و
میانگینش بالا می‌ماند.}
\end{table}

\paragraph{چرا این عدم‌تقارن ساختاری است.} سوال بعدی: آیا صندلی آمریکا
\emph{می‌تواند} حریف آینه‌ای را از پا در بیاورد؟ جواب نه است، و از خود
پارامترها درمی‌آید.

\begin{proposition}[صندلی آمریکا مسابقه‌ی فرسایش متقارن را می‌بازد]
\label{pr:asym}
فرض کنید حریف آینه‌ای باشد، یعنی $\sig_I=\sig_U=s$ و در نتیجه
$K_I=K_U=s/(1-\delta_K)=4s$. آنگاه برای \emph{هر}
$s\in\{0,0.25,0.5,0.75,1\}$، زمان فروپاشیِ گزاره~\ref{pr:collapse} در
نوع‌های میانگین برای ایران اکیداً بزرگ‌تر است.
\end{proposition}

\begin{proof}
هزینه‌ی خام در حالت پایدار
$R_i=c_{0i}+\kappa_i s+\alpha_i\cdot4s+\eta s^2$ است، پس
\[
  R_I = 1.4+6.2s+2.5s^2,\qquad R_U = 1.8+7.7s+2.5s^2 .
\]
با میانگین‌های پیشین $\mathbb{E}[\rbar_I]=5/7$، $\mathbb{E}[\rbar_U]=1/2$،
$\mathbb{E}[\mhat_I]=1/2$، $\mathbb{E}[\mhat_U]=3/5$ داریم $e_I^0=82.5$،
$e_U^0=76.5$، $\mbar_I=1.25$، $\mbar_U=1.44$، و از گزاره~\ref{pr:collapse}
با $x=1$: $T_i=0.675\,e_i^0\mbar_i/R_i$، یعنی صورت‌ها $69.6$ و $74.4$.
\end{proof}

\begin{table}[h]\centering\small
\begin{tabular}{lrrrrr}
\toprule
$s$ & $0$ & $0.25$ & $0.5$ & $0.75$ & $1.0$\\
\midrule
$T_I$ (مرحله) & $49.7$ & $22.4$ & $13.6$ & $9.3$ & $6.9$\\
$T_U$ (مرحله) & $41.3$ & $19.2$ & $11.8$ & $8.3$ & $6.2$\\
$T_I/T_U$     & $1.20$ & $1.17$ & $1.15$ & $1.13$ & $1.11$\\
\bottomrule
\end{tabular}
\caption{در برابر حریف آینه‌ای، آمریکا در هر سطح سیگنالی زودتر تمام می‌کند.
هرچه سیگنال بلندتر، برتری ایران کم‌تر — ولی هیچ‌وقت برنمی‌گردد. پس تسلیم
شدن صندلی آمریکا یک ضعف یادگیری نیست، \emph{یک واقعیت پارامتری} است.}
\end{table}

\paragraph{کانال \lr{(S4)}: تشخیص نوع متعهد.}
سه کانال باور بخش~\ref{sec:beliefs} همگی درباره‌ی $\theta_j$ حرف می‌زنند، و
دقیقاً به همین دلیل در برابر حریفی که یک \emph{قاعده‌ی پاسخ ثابت} دارد کور
می‌مانند: چنین قاعده‌ای مستقل از $\theta_j$ است، پس هیچ‌وقت تفکیک نمی‌شود.
آنچه او را لو می‌دهد خود قاعده است: $\sig_j(t)=\sig_i(t-1)$، مرحله به
مرحله. کانال چهارم یک \emph{پسین} روی همین فرضیه‌ی دودویی نگه می‌دارد — نه
یک شمارنده با آستانه. حریف آینه‌ای سیگنال قبلی ما را با احتمال $1$ تکرار
می‌کند و بقیه فقط تصادفی، پس هر تکرار دقیق شانس را در $\lambda=4$ ضرب
می‌کند و \emph{یک} عدم‌تطابق قاطعانه رد می‌کند:
\[
  \frac{P_{t+1}}{1-P_{t+1}}=\lambda\,\frac{P_t}{1-P_t}\quad
  \text{(تکرار دقیق)},\qquad P_{t+1}=0\quad\text{(هر چیز دیگر)},
\]
با پیشین $P_0=0.35$، یعنی سهم نوع‌های آینه‌ای در بین دو حریف مرجع منتشرشده.
دو ویژگی تازه به بردار باور اضافه شد (مجموعاً ۱۶): شماره‌ی ۱۴ همین
$P_t$، و شماره‌ی ۱۵ جایزه‌ای که خروج داوطلبانه تحویل حریف می‌دهد، یعنی
$\mathbb{E}_\mu[B_j]/175$ روی پسین نوع.

\paragraph{تلاش اول، و چرا شکست خورد.}
اولین جمله‌ای که نوشتیم دقیقاً همین بود: اگر اقدام \lr{EXIT} است \emph{و}
نوع متعهد شناسایی شده \emph{و} $\hat q$ پایین است، جریمه بگیر. سه شبکه با
آن آموزش دادیم. نتیجه:

\begin{table}[h]\centering\small
\begin{tabular}{lrrrr}
\toprule
چک‌پوینت & $J_I$ & $J_U$ & $d'$ & نرخ خروج داوطلبانه‌ی صندلی آمریکا\\
\midrule
پایه (بدون جمله) & $194.3$ & $-5.8$ & $156.8$ & $1.00$\\
با جمله، \lr{seed 11} & $187.0$ & $-12.8$ & $160.6$ & $1.00$\\
با جمله، \lr{seed 12} & $143.8$ & $-12.9$ & $160.5$ & $1.00$\\
\bottomrule
\end{tabular}
\caption{\lr{analysis/seat\_split.py}. نرخ خروج \emph{اصلاً تکان نخورد} و
$d'$ حتی \emph{بالا} رفت.}
\end{table}

\begin{findbox}
\textbf{یافته‌ی کلیدی ۶ — جریمه‌ای که می‌شود از آن فرار کرد، فقط فرار را
یاد می‌دهد.}
آن نسخه‌ی تشخیص‌دهنده به سه تکرار پشت سر هم نیاز داشت، یعنی زودترین شلیکش
$t=4$ بود. ولی عامل حدود $t\approx4$ خارج می‌شد. پس شبکه ساده‌ترین راه را
پیدا کرد: \emph{یک مرحله زودتر برو}. آن‌وقت جریمه هیچ‌وقت پرداخت نمی‌شود، و
چون حریف در $t$ کوچک‌تر تازه‌نفس‌تر است، جایزه‌اش $B_j$ هم \emph{بزرگ‌تر}
می‌شود. هدف دقیقاً برعکس شد. \textbf{درس کلی: هر آستانه‌ای که به یادگیرنده
یک مهلت بدهد، همان مهلت را یاد می‌دهد} — و برای همین در نسخه‌ی نهایی
آستانه را با پسین عوض کردیم، که مهلت ندارد و از همان $t=0$ روی
$P_0=0.35$ نشسته است.
\end{findbox}

\paragraph{تلاش دوم: به‌جای شرط‌گذاشتن، خودِ انتقال را قیمت‌گذاری کن.}
اشکال طرح اول این بود که جریمه یک \emph{دروازه‌ی زمانی} داشت. طرح دوم دروازه
ندارد: خروج داوطلبانه را به اندازه‌ی \emph{چیزی که تحویل می‌دهیم} جریمه
می‌کنیم،
\[
  -\ P_{\mathrm{concede}}\ \times\ \frac{\mathbb{E}_{\mu}\!\left[B_j\right]}{175},
  \qquad P_{\mathrm{concede}}=220 ,
\]
که در آن امید ریاضی روی همان پسین بخش~\ref{sec:beliefs} گرفته می‌شود
(ویژگی ۱۵). این جمله \emph{در مراحل اول بزرگ‌ترین مقدار} را دارد، چون حریف
آنجا تازه‌نفس است و $B_j$ نزدیک سقفش. یعنی «زودتر فرار کن» دیگر راه فرار
نیست، بدترین کار است.

جمله فقط روی \textbf{صندلی آمریکا} اعمال می‌شود. این قراردادی نیست:
گزاره~\ref{pr:asym} نشان می‌دهد آمریکا مسابقه‌ی متقارن را در هر سطح سیگنالی
می‌بازد، پس همان صندلی‌ای است که ساختار بازی به تسلیم هلش می‌دهد و تنها
جایی است که دو هدف از هم جدا می‌شوند.

\paragraph{و هزینه‌اش چقدر است؟}
قبل از آموزش، هزینه را حریف‌به‌حریف اندازه گرفتیم
(\lr{analysis/blanket\_cost.py}): چک‌پوینت پایه را برداشتیم و صندلی
آمریکایش را \emph{مجبور} کردیم خارج نشود.

\begin{table}[h]\centering\small
\begin{tabular}{lrrrr}
\toprule
حریف (صندلی ایران) & $D$ پایه & $D$ با نرفتن & $\Delta$ ما & $\Delta$ حریف\\
\midrule
\lr{Random}       & $108.1$ & $108.5$ & $+0.4$  & $-8.2$\\
\lr{TitForTat}    & $42.9$  & $-30.6$ & $\mathbf{-73.5}$ & $\mathbf{-192.4}$\\
\lr{MasterAgent}  & $36.1$  & $38.8$  & $+2.7$  & $-6.5$\\
\lr{refspec\_s3}  & $49.0$  & $49.7$  & $+0.7$  & $-103.4$\\
\lr{rich\_s0}     & $52.1$  & $57.1$  & $+5.0$  & $-10.6$\\
\lr{rich\_s1}     & $39.5$  & $46.6$  & $+7.1$  & $-66.3$\\
\lr{richshape\_s0}& $37.1$  & $28.4$  & $-8.6$  & $-40.6$\\
\bottomrule
\end{tabular}
\caption{در \emph{پنج حریف از هفت}، نرفتن حتی برای مطلوبیت خام خودمان هم
بهتر است. تنها جایی که تسلیم واقعاً به نفع ماست \lr{TitForTat} است — و
دقیقاً همان‌جاست که نباید تسلیم شویم.}
\end{table}

\begin{answer}
\textbf{این یک انتخاب است، نه یک کشف — ولی انتخاب گران‌قیمتی نیست.}
میانگین صندلی آمریکای ما $9.5$ واحد پایین می‌آید (‏یعنی امتیاز کل $4.7$)، و
تقریباً همه‌اش از یک رویارویی می‌آید. امتیاز \lr{TitForTat} در همان استخر
$13.7$ واحد پایین می‌آید. خالص: $+9.0$ به نفع ما.

و نکته‌ای که غافلگیرمان کرد: عامل پایه \emph{بیش از حد} تسلیم می‌شد. در پنج
مورد از هفت مورد، نرفتن هم برای رتبه بهتر است هم برای مطلوبیت خودمان.
گزاره‌ی ۱ (دروازه‌ی پایین) از اول همین را می‌گفت؛ خروج‌های زودهنگام عامل
یک اثر مرزیِ مهلت $\bar T=40$ بود، نه نتیجه‌ی تحلیل.
\end{answer}

\paragraph{و یک نشت که انتظارش را نداشتیم.}
جمله فقط روی صندلی آمریکا اعمال می‌شود، ولی وقتی شبکه‌ها را دوباره آموزش
دادیم صندلی \emph{ایران} هم عوض شد — و بدتر شد. علتش خودبازی است: صندلی
ایران در آموزش با همان صندلی آمریکای سرسخت‌شده‌ی خودمان بازی می‌کند، می‌بیند
که دیگر کوتاه نمی‌آید، و یاد می‌گیرد \emph{خودش} زودتر کوتاه بیاید.

\begin{table}[h]\centering\small
\begin{tabular}{lrrrr}
\toprule
سهم خودبازی $p_{\text{self}}$ & $J_I$ & نرخ خروج صندلی ایران & $b'$ & $J_U$\\
\midrule
$0.20$ (‏\lr{seed 23}) & $71.5$  & $0.73$ & $96.2$ & $97.4$\\
$0.10$ (‏\lr{seed 21}) & $79.7$  & $0.73$ & $93.9$ & $113.2$\\
$0.10$ (‏\lr{seed 22}) & $131.3$ & $0.48$ & $51.3$ & $108.0$\\
\bottomrule
\end{tabular}
\caption{هرچه خودبازی بیشتر، صندلی ایران تسلیم‌طلب‌تر. برای همین اجراهای
بعدی را با $p_{\text{self}}=0$ زدیم تا دو صندلی در آموزش کاملاً از هم جدا
بمانند — کاری که در این بازی بی‌خطر است، چون در ارزیابی هم هیچ‌وقت با خودمان
بازی نمی‌کنیم.}
\end{table}

\paragraph{تلاش سوم: جریمه‌ی بدون شرط، ارزیابی مرجع را می‌برد و تورنمنت را
می‌بازد.}
طرح دوم معیار مرجع را برد ($65.2$ در برابر $63.3$)، ولی وقتی همان چک‌پوینت
را در استخر بزرگ‌تر سنجیدیم معلوم شد گران تمام شده است:

\begin{table}[h]\centering\small
\begin{tabular}{lrrrr}
\toprule
 & \multicolumn{2}{c}{ارزیابی مرجع} & \multicolumn{2}{c}{استخر ۸ عاملی}\\
\cmidrule(lr){2-3}\cmidrule(lr){4-5}
چک‌پوینت & اختلاف & صندلی آمریکا & اختلاف & صندلی آمریکا\\
\midrule
نسخه‌ی ۱ (بدون جمله)     & $-10.8$ & $75.5$ & $-7.9$ & $52.1$\\
جمله‌ی بدون شرط (\lr{s22}) & $\mathbf{+1.9}$ & $36.6$ & $\mathbf{-16.6}$ & $31.9$\\
\bottomrule
\end{tabular}
\caption{جمله‌ی بدون شرط به \emph{هیچ} حریفی کوتاه نمی‌آید، نه فقط به آن یکی
که باید. صندلی آمریکا $20$ واحد از دست می‌دهد و در استخر کلاسی وضع از قبل
بدتر می‌شود.}
\end{table}

\begin{findbox}
\textbf{یافته‌ی کلیدی ۷ — «مجبورش کن نرود» با «یاد بگیرد نرود» یکی نیست.}
جدول \lr{blanket\_cost} بالا گفته بود نرفتن در پنج حریف از هفت حتی برای
مطلوبیت خودمان هم بهتر است. ولی آن آزمون سیاست \emph{موجود} را مجبور به
نرفتن می‌کرد؛ وقتی همان چیز را به یک شبکه \emph{یاد دادیم}، شبکه کل رفتارش
را عوض کرد: شروع کرد به جنگیدن، تاب‌آوری سوزاند، و نتیجه بدتر از اجبار ساده
شد. \textbf{اندازه‌گیری روی سیاست ثابت، پیش‌بینی سیاست آموزش‌دیده نیست.}
\end{findbox}

طرح نهایی هر دو اشکال را هم‌زمان حل می‌کند: جریمه را در پسینِ کانال
\lr{(S4)} ضرب می‌کنیم،
\[
  -\ P_{\mathrm{concede}}\ \times\ \underbrace{P_t}_{\text{ویژگی ۱۴}}
  \ \times\ \underbrace{\frac{\mathbb{E}_\mu[B_j]}{175}}_{\text{ویژگی ۱۵}} .
\]
حالا هیچ مهلتی برای فرار نیست (پسین از $t=0$ روی $0.35$ نشسته و جریمه‌ی
اولیه از سود خروج بزرگ‌تر است)، و در عین حال با \emph{اولین} پاسخ
غیرآینه‌ای پسین صفر می‌شود و جمله کاملاً خاموش. اندازه‌گیری مسیر پسین:

\paragraph{و یک شرط شناسایی که جا انداختنش همه‌چیز را خراب می‌کرد.}
آزمون «آیا حریف سیگنال قبلی ما را تکرار کرد؟» وقتی آن سیگنال \emph{صفر}
باشد \emph{بی‌معنی} است: حریف آرامی که همیشه $\sig=0$ می‌زند، بی‌آنکه آینه
باشد، صفرِ ما را تکرار می‌کند. اولین نسخه‌ی پسین همین اشتباه را داشت و
نتیجه‌اش این شد که صندلی آمریکا به \emph{هیچ‌کس} کوتاه نمی‌آمد. شرط درست
این است که مرحله فقط وقتی شاهد حساب شود که \emph{چیزی برای کپی کردن} وجود
داشته باشد، یعنی $\sig_i(t-1)>0$. اثرش:

\begin{table}[h]\centering\small
\begin{tabular}{lrrrrrr}
\toprule
حریف & $t{=}0$ & $1$ & $2$ & $3$ & $4$ & $5$\\
\midrule
\lr{TitForTat}    & $0.35$ & $0.35$ & $0.68$ & $0.90$ & $0.97$ & $\mathbf{0.99}$\\
\lr{Random}       & $0.35$ & $0.35$ & $0.68$ & $\mathbf{0.00}$ & --- & ---\\
\lr{refspec\_s3}  & $0.35$ & $0.35$ & $0.68$ & $0.90$ & $\mathbf{0.00}$ & $0.00$\\
\lr{rich\_s0}     & $0.35$ & $0.35$ & $\mathbf{0.00}$ & $0.00$ & $0.00$ & $0.00$\\
\lr{MasterAgent}  & $0.35$ & $0.35$ & $\mathbf{0.00}$ & $0.00$ & $0.00$ & $0.00$\\
\bottomrule
\end{tabular}
\caption{$P_t$، پسین «حریف قاعده‌ی آینه‌ای دارد». فقط \lr{TitForTat} به یک
می‌رود (جریمه از $-73$ به $-195$)؛ هر پنج حریف دیگر — از جمله سه شبکه‌ی
آموزش‌دیده و بنچمارک استاد — با اولین پاسخ ناسازگار به صفر می‌افتند و
جمله کاملاً خاموش می‌شود.}
\end{table}

\paragraph{انتخاب روی بذرهای دیده‌نشده.}
دستورالعمل (بند ۵-۴) می‌گوید اجرای نهایی «از بذرهای تازه و افشا نشده استفاده
خواهد کرد». پس برتری‌ای که فقط روی یک فهرست بذر مشخص وجود داشته باشد
بی‌ارزش است. با \lr{analysis/seed\_robust.py} هر نامزد را روی شش بلوک بذر
\emph{مجزا} — که پنج‌تای‌شان را اصلاً در تنظیم ندیده بودیم — دوباره سنجیدیم:

\begin{table}[h]\centering\small
\begin{tabular}{lrrrrrr}
\toprule
چک‌پوینت & $s{=}7$ & $10^5$ & $107919$ & $115838$ & $123757$ & $131676$\\
\midrule
نسخه‌ی ۱ (‏\lr{refspecfull})       & $-10.8$ & $-11.3$ & $-8.8$ & $-8.0$ & $-12.5$ & $-16.4$\\
\lr{hold\_s21} ($p_{\text{self}}{=}0.10$) & $-9.7$ & $-5.3$ & $-6.8$ & $-16.5$ & $-3.3$ & $-3.0$\\
\lr{hold\_s23} ($p_{\text{self}}{=}0.20$) & $-15.7$ & $-15.5$ & $-11.5$ & $-21.5$ & $-8.0$ & $-11.7$\\
\lr{hold\_s22} ($p_{\text{self}}{=}0.10$) & $+1.9$ & $+7.9$ & $+6.7$ & $-3.1$ & $+4.3$ & $+2.5$\\
\bottomrule
\end{tabular}
\caption{اختلاف با \lr{TitForTat} روی شش بلوک بذر مجزا. نسخه‌ی ۱ در
\emph{هر شش} بلوک می‌بازد؛ \lr{hold\_s22} در پنج بلوک از شش می‌برد
(میانگین $+3.4$). یعنی برتری واقعی است ولی هنوز حاشیه‌اش کم است — و دقیقاً
به همین دلیل سراغ ترمیم صندلی ایران رفتیم.}
\end{table}

\begin{findbox}
\textbf{یک نتیجه‌ی منفی — و اینکه چطور برطرف شد.}
در رژیم اول (استخر استاندارد، $300$ هزار گام) ویژگی‌های باور — که از نظر
آماری \emph{کار می‌کنند} و خطای برآورد را $72.9\%$ کم می‌کنند — عملکرد را
\textbf{بهتر نکردند}: «فقط شکل‌دهی» $90.6$ گرفت و «ویژگی + شکل‌دهی»
$83.2$.

فرضیه‌ی ما این بود که مشکل \emph{کیفیت باور} نیست بلکه \emph{بودجه‌ی نمونه}
است: شبکه یک \lr{GRU} است و خودش پسین \emph{ضمنی} حمل می‌کند، و ۱۴ ویژگی
اضافه ورودی را از ۱۶ به ۳۰ بعد می‌برد.

این فرضیه \emph{آزمون‌پذیر} بود و آزمونش را پاس کرد: با $700$ هزار گام و
استخر هماهنگ، «ویژگی + شکل‌دهی» به $89.8$ رسید و از «فقط شکل‌دهی» جلو زد
(‏$84.8$ در برابر $83.8$ در تنظیم رسمی، و مهم‌تر اینکه امتیاز
\lr{TitForTat} را از $109.8$ به $95.9$ پایین آورد). یعنی ویژگی‌های باور
با بودجه‌ی کافی \emph{می‌ارزند} — و به همین دلیل چک‌پوینت ارسالی نظریه‌ی
کامل است، نه نسخه‌ی حذفی.
\end{findbox}

\subsection{تلاش چهارم: اگر جریمه شرطی است، چرا فقط یک صندلی؟}
\label{sec:both}

جمله‌ی نهایی بعد از اصلاح شرط شناسایی، مقابل حریفی که آینه نیست \emph{اصلاً
روشن نمی‌شود}. ولی محدود کردنش به صندلی آمریکا از دورانی مانده بود که جریمه
بدون شرط بود و باید هزینه‌اش را جایی مهار می‌کردیم. حالا که هزینه‌ی بی‌مورد
ندارد، آن محدودیت دلیلش را از دست داده است: هر دو صندلی وقتی به یک نوع
متعهدِ شناسایی‌شده کوتاه می‌آیند \emph{همان} جایزه را تحویل می‌دهند — و آن
جایزه دقیقاً همان $b'$ در تعریف $J_I$ است که تا اینجا هیچ‌وقت جبرانش نکرده
بودیم. پس \lr{CONCEDE\_SEAT} را برداشتیم و جمله را روی هر دو صندلی گذاشتیم.

\begin{table}[h]\centering\small
\begin{tabular}{lrrrl}
\toprule
نسخه & $J_I$ & $J_U$ & مجموع & در برابر آستانه‌ی $231.8$\\
\midrule
پایه (بدون جمله)               & $194.3$ & $-5.8$  & $188.5$ & می‌بازد\\
پسین بی‌شرط، فقط صندلی آمریکا   & $170.8$ & $93.7$  & $264.5$ & می‌برد\\
پسین شرطی، فقط صندلی آمریکا     & $123.6$ & $81.6$  & $205.2$ & می‌بازد\\
\textbf{پسین شرطی، هر دو صندلی} & $\mathbf{189.3}$ & $\mathbf{84.1}$ & $\mathbf{273.4}$ & می‌برد\\
\bottomrule
\end{tabular}
\caption{بهترین صندلی هر نسخه از میان بذرهای آموزش‌دیده
(\lr{analysis/seat\_split.py}، $80$ اپیزود). ستون $J_I$ همان چیزی است که
تلاش چهارم هدف گرفت و از $123.6$ به $189.3$ رفت — با $b'=-23.8$، یعنی
\lr{TitForTat} در برابر صندلی ایران ما مطلوبیت \emph{منفی} می‌گیرد. و
برخلاف نسخه‌ی بی‌شرط، $J_U$ هم بالا ماند؛ یعنی این بار هر دو صندلی
هم‌زمان درست شدند.}
\end{table}

\paragraph{نتیجه‌ی رسمی.}
چک‌پوینت نهایی از چسباندن بهترین صندلی ایران و بهترین صندلی آمریکا ساخته
شده (\lr{analysis/splice.py}؛ اسکریپت اگر سورس نظریه‌ی دو طرف بایت‌به‌بایت
یکی نباشد کار نمی‌کند). خروجی خودِ \lr{scripts/evaluate.py}:

\begin{table}[h]\centering
\begin{tabular}{lrrr}
\toprule
رتبه & امتیاز & به‌عنوان ایران & به‌عنوان آمریکا\\
\midrule
\textbf{چک‌پوینت ارسالی} & $\mathbf{71.8}$ & $82.6$ & $\mathbf{61.0}$\\
\lr{TitForTat} (مرجع)    & $61.1$ & $83.1$ & $39.1$\\
\lr{Random} (مرجع)       & $-7.8$ & $-10.0$ & $-5.6$\\
\bottomrule
\end{tabular}
\caption{معیار «امتیاز بالاتر از \lr{TitForTat}» با اختلاف $+10.7$ برقرار
است. برتری تقریباً همه‌اش از صندلی \emph{آمریکا} می‌آید ($61.0$ در برابر
$39.1$) — همان صندلی‌ای که در ابتدای کار $100\%$ تسلیم می‌شد.}
\end{table}

و در استخر چهارتایی که بنچمارک استاد هم در آن هست
(\lr{analysis/benchmark\_table.py}، $60$ اپیزود):

\begin{table}[h]\centering\small
\begin{tabular}{lrrrl}
\toprule
عامل & امتیاز & ایران & آمریکا & بازه‌ی اطمینان $95\%$\\
\midrule
\textbf{چک‌پوینت ارسالی} & $\mathbf{65.2}$ & $76.4$ & $54.0$ & $[56.5,\ 72.7]$\\
\lr{TitForTat} (مرجع)    & $64.2$ & $83.4$ & $45.1$ & $[54.7,\ 74.3]$\\
\lr{MasterAgent} (بنچمارک استاد) & $35.0$ & $41.9$ & $28.0$ & $[27.8,\ 42.7]$\\
\lr{Random} (مرجع)       & $4.5$  & $-0.5$ & $9.5$  & $[-1.0,\ 10.1]$\\
\bottomrule
\end{tabular}
\caption{بنچمارک استاد را در \emph{هر دو} نقش می‌بریم ($63.8$ در برابر
$-6.4$ وقتی ما ایرانیم، و $40.1$ در برابر $9.0$ وقتی ما آمریکاییم) و
بازه‌های اطمینانمان هم‌پوشانی ندارند. رویارویی با \lr{TitForTat} وقتی ما
آمریکاییم هنوز تنها جایی است که عقبیم، ولی فاصله‌اش از $115$ واحد در
نسخه‌ی ۱ به $25.8$ واحد رسیده است.}
\end{table}

\paragraph{و آیا این برتری روی بذرهای دیده‌نشده هم هست؟}
همان آزمون شش‌بلوکی، این بار روی هر سه نامزد نهایی:

\begin{table}[h]\centering\small
\begin{tabular}{lrrrrrrr}
\toprule
چک‌پوینت & $s{=}7$ & $10^5$ & $107919$ & $115838$ & $123757$ & $131676$ & میانگین\\
\midrule
پسین بی‌شرط (یک صندلی) & $+10.2$ & $+6.0$ & $+12.4$ & $-1.0$ & $+10.9$ & $+11.8$ & $+8.4$\\
شرطی هر دو صندلی، دسته‌ی اول & $+3.2$ & $+2.4$ & $+6.5$ & $-0.1$ & $+5.1$ & $+1.5$ & $+3.1$\\
\textbf{ارسالی} & $\mathbf{+10.7}$ & $+11.2$ & $+17.0$ & $+8.4$ & $+14.2$ & $+8.9$ & $\mathbf{+11.7}$\\
\bottomrule
\end{tabular}
\caption{چک‌پوینت ارسالی در \emph{هر شش} بلوک می‌برد و بدترین بلوکش هم
$+8.4$ است؛ پنج بلوک از این شش‌تا را اصلاً در تنظیم ندیده بودیم. مقایسه‌ی دو
ردیف اول نشان می‌دهد پراکندگی بین بذرها در این بازی زیاد است — و دقیقاً به
همین دلیل انتخاب چک‌پوینت را روی همین آزمون انجام دادیم، نه روی یک اجرای
واحد.}
\end{table}

\begin{findbox}
\textbf{یافته‌ی کلیدی ۸ — یک معیار کافی نیست.}
ارزیابی مرجع فقط سه عامل دارد. تورنمنت کلاسی ده‌ها. همان چک‌پوینت‌ها را در
استخرهای بزرگ‌تر هم سنجیدیم (\lr{analysis/classpool.py}؛ چک‌پوینت‌های
دیگرمان و بنچمارک استاد به‌جای هم‌کلاسی‌ها):

\vspace{4pt}
\centering\small
\begin{tabular}{lrrrrrr}
\toprule
اندازه‌ی استخر & $3$ & $4$ & $5$ & $6$ & $7$ & $8$\\
\midrule
پسین بی‌شرط (یک صندلی) & $+3.4$ & $-6.4$ & $-17.2$ & $-19.6$ & $-21.4$ & ---\\
شرطی هر دو صندلی، دسته‌ی اول & $+2.9$ & $-4.9$ & $-9.7$ & $-9.2$ & $-8.3$ & $-7.8$\\
\textbf{ارسالی} & $\mathbf{+9.9}$ & $-1.1$ & $-4.0$ & $-3.3$ & $-3.8$ & $\mathbf{-4.0}$\\
\bottomrule
\end{tabular}

\vspace{5pt}
\raggedright
نسخه‌ی بی‌شرط هرچه استخر بزرگ‌تر می‌شود \emph{بدتر} می‌شود، چون به هیچ‌کس
کوتاه نمی‌آید و مقابل حریفانی که اصلاً نمی‌شود از پا درشان آورد تاب‌آوری
می‌سوزاند: صندلی آمریکایش از $44$ به $11$ سقوط می‌کند. چک‌پوینت ارسالی
صندلی آمریکایش را روی $43$ تا $49$ نگه می‌دارد و اختلافش با \lr{TitForTat}
حدود $4$ واحد تثبیت می‌شود.
\end{findbox}

\begin{answer}
\textbf{چه چیزی این تفاوت را می‌سازد.} همان \emph{شرطی} بودن جریمه.
\begin{enumerate}[leftmargin=1.6em,itemsep=2pt]
  \item جریمه‌ی بی‌شرط دقیقاً همان چیزی است که بند ۵-۲ دستورالعمل درباره‌اش
        هشدار می‌دهد: راهنمایی‌ای که «صرفاً برای بقا» پاداش می‌دهد. عامل
        یاد می‌گیرد به \emph{هیچ‌کس} کوتاه نیاید، حتی به حریفی که هیچ
        دلیلی برای محروم کردنش وجود ندارد و هیچ شانسی هم برای بردنش.
  \item با ضرب کردن در پسین، جمله فقط مقابل نوع متعهدِ شناسایی‌شده روشن
        می‌ماند. جدول $P_t$ بالا نشان می‌دهد این تفکیک واقعاً کار می‌کند:
        فقط \lr{TitForTat} به یک می‌رود.
  \item و برداشتن قید «فقط صندلی آمریکا» همان کسری $J_I$ را جبران کرد که
        دو تلاش قبلی نتوانسته بودند.
\end{enumerate}
هر سه معیاری که داشتیم — ارزیابی مرجع، بذرهای دیده‌نشده، و استخر کلاسی —
یک برنده دارند، و آن همین نسخه است. برای همین دیگر بین دو هدف مجبور به
انتخاب نشدیم.
\end{answer}

\begin{answer}
\textbf{یک نکته درباره‌ی روش انتخاب.} در تمام این کار، انتخاب بین
چک‌پوینت‌ها را روی \emph{بذرهای دیده‌نشده} انجام دادیم، نه روی اجرایی که
بهترین عدد را داده بود. مقایسه‌ی دو دسته‌ی «شرطی، هر دو صندلی» در جدول
بالا نشان می‌دهد چرا: دو دسته با \emph{کد یکسان} و فقط بذرهای متفاوت،
میانگین $+3.1$ و $+11.7$ گرفتند. اگر فقط یک دسته را می‌زدیم و همان را
گزارش می‌کردیم، عدد گزارش‌شده بیشتر درباره‌ی شانس بود تا درباره‌ی نظریه.
\end{answer}

\subsection{تشخیص چهارم: معیاری که تمام مدت اشتباه اندازه می‌گرفتیم}
\label{sec:realboard}

تا اینجای گزارش، هر عددی که گزارش کرده‌ایم با \lr{scripts/evaluate.py}
گرفته شده — همان دستوری که بند ۵-۴ دستورالعمل می‌گوید قبل از ارسال اجرا
کنید. بعد چک‌پوینت را روی \lr{seeboard.hpc.ipm.ac.ir} بارگذاری کردیم و
برد چیز دیگری گفت:

\begin{table}[h]\centering
\begin{tabular}{lrr}
\toprule
 & \lr{evaluate.py} (سه عاملی) & برد کلاسی\\
\midrule
اختلاف با \lr{TitForTat} & $\mathbf{+10.7}$ & $\mathbf{-0.9}$\\
\bottomrule
\end{tabular}
\caption{یک چک‌پوینت، دو معیار، دو علامت مخالف.}
\end{table}

علتش را برد خودش نوشته بود: «امتیاز = میانگین مطلوبیت خام در برابر
\emph{استخر مرجع} (پنج حریف اسکریپتی به‌علاوه‌ی بنچمارک استاد)». ولی
\lr{see/agents/scripted.py} فقط \emph{دو} حریف دارد. سه حریف دیگر
(\lr{Hawk}، \lr{Dove}، \lr{BayesianThreshold}) در بند ۱۱ سند
\lr{DESIGN.md} توصیف شده‌اند و همان‌جا نوشته که از مخزن حذف شده‌اند. ما آن
جمله را یک نکته‌ی تاریخی خوانده بودیم، نه ترکیب استخر ارزیابی.

\begin{findbox}
\textbf{یافته‌ی کلیدی ۹ — بهینه‌سازی روی معیار در دسترس، نه معیار واقعی.}
تمام تنظیم بخش~\ref{sec:tourn} روی استخر سه‌تایی انجام شده بود. آن استخر
معیار نمره‌دهی نیست، و بدتر اینکه در این مسئله با آن \emph{هم‌جهت هم
نیست}. این جدی‌ترین اشتباه روش‌شناختی این پروژه بود و بقیه‌ی این بخش
تصحیح آن است.
\end{findbox}

\paragraph{بازسازی استخر مرجع.}
سه حریف گم‌شده را از روی توصیف \lr{DESIGN.md} نوشتیم
(\lr{my\_team/reference\_pool.py}) و پارامترهای آزادشان را روی شش سطری
که برد منتشر کرده بود برازش دادیم (\lr{analysis/board.py}). برازش یک
چیز ساختاری را هم مشخص کرد که هیچ‌جا نوشته نشده بود:

\begin{keybox}
\textbf{امتیاز هر حریف مرجع، شامل بازی با \emph{خودش} هم هست.}
بدون این فرض \lr{Hawk} با اختلاف زیاد اول می‌شود؛ با آن سوم می‌شود،
دقیقاً مثل برد. دلیلش همان جمله‌ی \lr{DESIGN.md} است که «عامل
همیشه‌تشدیدکننده در بازی آینه‌ای خودش را نابود می‌کند». و چون \emph{ما}
عضو استخر مرجع نیستیم، هیچ‌وقت با خودمان بازی نمی‌کنیم.
\end{keybox}

بعد برد جدول رویارویی‌های خودِ چک‌پوینت ما را هم منتشر کرد، که دوازده عدد
دقیق برای یک چک‌پوینت \emph{ثابت} می‌دهد — یک مسئله‌ی برازش بسیار
خوش‌وضع‌تر (\lr{analysis/fit\_refs.py}). نتیجه:

\begin{table}[h]\centering\small
\begin{tabular}{lrrl}
\toprule
حریف مرجع & برد واقعی & بازسازی ما & وضعیت\\
\midrule
\lr{Random}            & $119.8$ / $113.8$ & $119.8$ / $113.8$ & \textbf{دقیق} (در مخزن هست)\\
\lr{TitForTat}         & $34.3$ / $2.7$    & $34.3$ / $2.7$    & \textbf{دقیق} (در مخزن هست)\\
\lr{MasterAgent}       & $70.1$ / $-2.7$   & $70.1$ / $-2.7$   & \textbf{دقیق} (بازیابی‌شده)\\
\lr{Hawk}              & $-12.1$ / $-58.7$ & $-12.0$ / $-58.7$ & \textbf{دقیق} (همیشه $\sig=1$)\\
\lr{Dove}              & $144.4$ / $135.9$ & $154.8$ / $130.6$ & تقریبی\\
\lr{BayesianThreshold} & $77.6$ / $9.4$    & $82.1$ / $35.6$   & تقریبی\\
\bottomrule
\end{tabular}
\caption{چهار حریف از شش \emph{دقیقاً} بازتولید می‌شوند. سه‌تای اول را از
قبل دقیق داشتیم، پس تطابقشان بذرها و تعداد اپیزود و پروتکل را هم تایید
می‌کند؛ \lr{Hawk} با ساده‌ترین تعریف ممکن درست درآمد. کل بازسازی حدود
$+3$ خوش‌بین است، پس برای انتخاب چک‌پوینت حاشیه‌ی راحت لازم داریم.}
\end{table}

\paragraph{تشخیص، از روی همان جدول.}
کل کسری در یک رویارویی بود: در صندلی آمریکا مقابل \lr{Hawk} امتیاز
$-58.7$ می‌گرفتیم، در حالی که صرفاً گرفتن گزینه‌ی خروج آنجا حدود $+20$
می‌ارزد. یعنی \emph{یک} حریف حدود $13$ واحد از میانگین ما می‌خورد، و کل
فاصله تا \lr{TitForTat} کمتر از یک واحد بود. علتش هم روشن بود: استخر
آموزش ما \lr{Random} و \lr{TitForTat} بود و عامل ما \textbf{هیچ‌وقت با
حریفی که حاضر نیست برود تمرین نکرده بود}.

\paragraph{اصلاح اول: جریمه باید انتقال \emph{مورد انتظار} باشد.}
جمله‌ی جریمه‌ی تسلیم، خروج داوطلبانه را به اندازه‌ی $B_j$ جریمه می‌کرد.
ولی تسلیم شدن فقط وقتی چیزی منتقل می‌کند که ما قرار بوده زنده بمانیم و
جایزه را بگیریم؛ مقابل حریفی که نمی‌توانیم از پا در بیاوریم، او آن جایزه
را در هر صورت می‌گیرد. پس انتقال مورد انتظار $\hat q\cdot B_j$ است نه
$B_j$. ضریب $\hat q$ (که از قبل ویژگی شماره ۸ بود) اضافه شد.

\paragraph{اصلاح دوم، و مهم‌تر: معیار اصلاً نسبی نیست.}
اصلاح اول جواب داد ولی سوال بزرگ‌تری را باز کرد. امتیاز هر حریف مرجع از
بازی‌هایش با \emph{استخر مرجع} می‌آید، و ما عضو آن استخر نیستیم. پس
بازی ما با \lr{TitForTat} اصلاً در امتیاز او حساب نمی‌شود:

\begin{keybox}
\textbf{امتیاز \lr{TitForTat} یک عدد ثابت است ($53.8$) که هیچ کاری از ما
تغییرش نمی‌دهد.} معیار «امتیاز بالاتر از \lr{TitForTat}» پس به این
فرو‌ می‌آید:
\[
  \text{میانگین مطلوبیت \emph{خودمان} روی استخر مرجع}\ >\ 53.8 .
\]
هیچ جمله‌ی محروم‌سازی‌ای در کار نیست.
\end{keybox}

این دقیقاً پایه‌ی جمله‌ی دوم شکل‌دهی را می‌زند. آن جمله برای معیاری نوشته
شده بود که مطلوبیت ما را با مطلوبیت حریف مقایسه می‌کند
(رابطه‌ی~\eqref{eq:crit})، و آن رابطه برای استخر سه‌تایی
\lr{evaluate.py} \emph{درست} است — آنجا امتیاز \lr{TitForTat} شامل بازی
با ما هم می‌شود. برای برد کلاسی غلط است. روی برد، آن جمله فقط مطلوبیت
خودمان را می‌فروخت بدون اینکه چیزی بخرد.

پس قیمتش را صفر گذاشتیم. نتیجه‌ی جانبی‌اش این است که شکل‌دهی پاداش دوباره
\emph{کاملاً پتانسیل‌پایه} می‌شود و طبق قضیه‌ی
\lr{Ng, Harada \& Russell} سیاست بهینه را ثابت نگه می‌دارد — یعنی عامل
دقیقاً همان کمیتی را بیشینه می‌کند که برد می‌سنجد. ثابت
\lr{CONCEDE\_PRICE} را حذف نکردیم چون برای یک قاعده‌ی \emph{واقعاً} نسبی
درست است؛ تورنمنت دوره‌ای کلاسی یکی از آن‌هاست، هرچند وزنش آنجا فقط
$1/(N-1)$ است.

\paragraph{نتیجه.}
دو تغییر: استخر آموزش به همان شش حریف برد عوض شد
(\lr{--pool board})، و جمله‌ی جریمه صفر شد.

\begin{table}[h]\centering\small
\begin{tabular}{lrrrrr}
\toprule
 & \multicolumn{4}{c}{امتیاز روی بازسازی، به تفکیک پایه‌ی بذر} & \\
\cmidrule(lr){2-5}
چک‌پوینت & $20260901$ & $7\!\times\!10^6$ & $31415926$ & $88888888$ & میانگین\\
\midrule
آپلود اول (استخر سه‌تایی) & $54.1$ & $58.1$ & $53.4$ & $59.2$ & $56.2$\\
استخر برد + جریمه‌ی $\hat q$-دار & $68.1$ & $67.5$ & $65.8$ & $70.2$ & $67.9$\\
\textbf{استخر برد + بدون جریمه} & $\mathbf{68.8}$ & $\mathbf{69.8}$ & $\mathbf{69.4}$ & $\mathbf{72.3}$ & $\mathbf{70.1}$\\
\bottomrule
\end{tabular}
\caption{چهار پایه‌ی بذر مجزا؛ پایه‌ی $20260901$ همان بذر تمرینی
اعلام‌شده است و سه‌تای دیگر را خودمان انتخاب کردیم تا مطمئن شویم برتری
به بذر خاصی وابسته نیست. آپلود اول روی برد واقعی $52.9$ گرفت و بازسازی
$54.1$ می‌گفت.}
\end{table}

\begin{answer}
\textbf{پیش‌بینی ما برای این چک‌پوینت $\approx67$ بود. برد $56.8$ داد.}
آپلود دوم را انجام دادیم و عدد واقعی آمد؛ پیش‌بینی حدود ده واحد خوش‌بین
بود. جدول بعد نشان می‌دهد کل این خطا کجا نشسته است، و آن جدول از هر
پیش‌بینی‌ای که در این گزارش نوشته‌ایم آموزنده‌تر است.

یک هشدار هم سر جایش می‌ماند: چک‌پوینت نهایی روی \lr{scripts/evaluate.py}
— همان استخر سه‌تایی — \emph{بدتر} از قبلی است ($65.7$ در برابر $86.8$).
این تناقض نیست، خودِ یافته است: آن دو معیار در این مسئله در جهت مخالف هم
حرکت می‌کنند.
\end{answer}

\subsection{خواندن دوم برد: کجای پیش‌بینی غلط بود}
\label{sec:board2}

آپلود دوم هم امتیاز کل داد و هم جدول به‌تفکیک‌حریف. چون این جدول برای یک
چک‌پوینت \emph{ثابت} است، دوازده عدد دقیق به ما می‌دهد و می‌شود بازسازی
را نقطه‌به‌نقطه محاکمه کرد.

\begin{table}[h]\centering\small
\begin{tabular}{lrrrrrr}
\toprule
حریف مرجع & \multicolumn{2}{c}{آپلود اول (واقعی)}
           & \multicolumn{2}{c}{آپلود دوم (واقعی)}
           & \multicolumn{2}{c}{بازسازی برای آپلود دوم}\\
\cmidrule(lr){2-3}\cmidrule(lr){4-5}\cmidrule(lr){6-7}
 & ایران & آمریکا & ایران & آمریکا & ایران & آمریکا\\
\midrule
\lr{Random}            & $119.8$ & $113.8$ & $121.4$ & $113.3$ & $121.4$ & $113.3$\\
\lr{TitForTat}         & $34.3$  & $2.7$   & $34.1$  & $9.7$   & $34.1$  & $9.7$\\
\lr{Hawk}              & $-12.1$ & $-58.7$ & $12.0$  & $4.7$   & $12.0$  & $4.7$\\
\lr{MasterAgent}       & $70.1$  & $-2.7$  & $96.8$  & $44.4$  & $96.8$  & $44.4$\\
\midrule
\lr{Dove}              & $144.4$ & $135.9$ & $104.9$ & $64.6$  & $\mathbf{152.8}$ & $\mathbf{143.3}$\\
\lr{BayesianThreshold} & $77.6$  & $9.4$   & $60.9$  & $15.0$  & $\mathbf{76.4}$  & $\mathbf{34.3}$\\
\midrule
\textbf{میانگین}       & $72.3$  & $33.4$  & $71.7$  & $42.0$  & $82.3$ & $58.3$\\
\textbf{امتیاز}        & \multicolumn{2}{c}{$52.9$}
                       & \multicolumn{2}{c}{$\mathbf{56.8}$}
                       & \multicolumn{2}{c}{$70.3$}\\
\bottomrule
\end{tabular}
\caption{چهار حریف بالای خط \emph{تا رقم اعشار} بازتولید می‌شوند — نه
تقریباً، دقیقاً. کل خطای $13.5$ واحدی پیش‌بینی در دو ردیف پایین نشسته
است.}
\label{tab:board2}
\end{table}

این نتیجه دو رو دارد. روی خوبش: چهار حریف از شش دقیقاً درست‌اند، پس
پروتکل و بذرها و شمار اپیزودها همه درست بازسازی شده‌اند و آن بخش از
تخمین قابل اتکاست. روی بدش: پیش‌بینی کل، بی‌فایده بود، چون دو حریف
باقی‌مانده تمام خطا را می‌سازند.

\paragraph{چرا برازش دوباره جواب نداد.}
با دو آپلود، حالا برای \lr{Dove} چهار عدد هدف داریم نه دو تا. ولی
\lr{analysis/fit\_refs2.py} نشان داد هیچ مقداری از $\theta$ هر دو را
نمی‌سازد: \lr{Dove} واقعی به دو آپلود ما امتیازهایی با فاصله‌ی
$-39.5/-71.3$ می‌دهد، در حالی که بازسازی ما در هر مقدار $\theta$ فاصله را
زیر $9$ نگه می‌دارد. یعنی چیزی که کم داریم یک \emph{پارامتر} نیست، یک
\emph{سازوکار} است.

\paragraph{سازوکار گمشده، سمت خودمان بود.}
به‌جای حدس زدن دوباره‌ی حریف، رفتار خودمان را مقابل یک حریف کاملاً منفعل
اندازه گرفتیم (\lr{analysis/sigprofile.py}):

\begin{table}[h]\centering\small
\begin{tabular}{lrr}
\toprule
آماره & آپلود اول & آپلود دوم\\
\midrule
میانگین سیگنال                & $0.06$ & $\mathbf{0.75}$\\
کسر مراحل صرف‌شده روی \lr{HOLD} & $0.24$ & $\mathbf{0.96}$\\
کسر اپیزودهایی که خودمان اول خارج می‌شویم & $0.24$ & $\mathbf{0.00}$\\
\bottomrule
\end{tabular}
\caption{چک‌پوینت ارسالی $\sigma=0.75$ را در مرحله‌ی صفر می‌گذارد، ۹۶٪
مراحل نگهش می‌دارد و هیچ‌وقت خودش خارج نمی‌شود.}
\end{table}

مقابل حریفی که زود می‌رود این رفتار مجانی است. مقابل حریفی که فقط
\emph{صبر می‌کند}، تمام اپیزود هزینه‌ی سیگنال می‌دهیم و گزینه‌ی خروج را
هم دور می‌ریزیم. هر \lr{Dove}ای که ما نوشتیم بلافاصله می‌رود، برای همین
بازسازی اصلاً نمی‌توانست این تفاوت را ببیند.

\begin{findbox}
\textbf{یافته‌ی کلیدی ۱۰ — استخر را \emph{جایگزین} کردیم به‌جای اینکه
\emph{گسترش} بدهیم.} وقتی مربی را به \lr{--pool board} بردیم، کل دیکشنری
حریف‌ها عوض شد و دو کهن‌الگوی \lr{Patient} و \lr{Plateau} — تنها
حریف‌های منفعلِ صبور — بی‌سروصدا از استخر افتادند. عامل حاصل یاد گرفت که
تسلیم شدن هزینه‌ای ندارد، چون دیگر هیچ‌کس نبود که فقط بایستد و صبر کند.
بهای این حذف روی برد، حدود $13$ واحد بود. درسِ عملی: این دسته خطا را با
\emph{امتیاز} پیدا نمی‌کنید، با اندازه‌گیری \emph{رفتار} پیدا می‌کنید.
\end{findbox}

\paragraph{وقتی بازسازی جواب نمی‌دهد، دنبال جانشینِ هم‌ترتیب بگردید.}
\lr{analysis/proxy\_probe.py} همه‌ی حریف‌های محلی را امتحان کرد تا ببیند
کدام‌شان دو آپلود ما را \emph{به همان ترتیبِ} برد می‌چیند. فقط
\lr{Patient} این کار را می‌کند (فاصله‌ی $-18.6/-43.3$ در برابر
$-39.5/-71.3$ برد)؛ همه‌ی نسخه‌های \lr{Dove} فاصله‌ی تقریباً صفر می‌دهند.
سطح مطلق این حریف اصلاً شبیه \lr{Dove} نیست و ادعایی هم نمی‌کنیم — ولی
برای \emph{انتخاب بین چک‌پوینت‌ها} ترتیب کافی است، و ترتیب همان چیزی بود
که نداشتیم.

بر همین اساس معیار انتخاب را عوض کردیم (\lr{analysis/select2.py}): دیگر
شواهد قابل‌اعتماد و غیرقابل‌اعتماد را با هم میانگین نمی‌گیریم. یک ستون
برای چهار حریف دقیق، یک ستون برای جانشین صبور، و شرط آپلود این است که
کاندیدا روی \emph{هر دو} بهتر باشد. دقیقاً همین معامله‌ی یک‌طرفه بود که
چک‌پوینت فعلی انجام داد و برد بابتش جریمه گرفت.

\subsection{خواندن سوم: تصحیح جواب داد، و یک سقف پیدا شد}
\label{sec:board3}

استخر آموزش را با کهن‌الگوهای صبور گسترش دادیم (\lr{--pool board+}) و
چک‌پوینت حاصل را آپلود کردیم. نتیجه سه چیز را نشان داد.

\begin{table}[h]\centering\small
\begin{tabular}{lrrrrrr}
\toprule
حریف مرجع & \multicolumn{2}{c}{آپلود اول}
           & \multicolumn{2}{c}{آپلود دوم}
           & \multicolumn{2}{c}{آپلود سوم}\\
\cmidrule(lr){2-3}\cmidrule(lr){4-5}\cmidrule(lr){6-7}
 & ایران & آمریکا & ایران & آمریکا & ایران & آمریکا\\
\midrule
\lr{Random}            & $119.8$ & $113.8$ & $121.4$ & $113.3$ & $116.5$ & $62.1$\\
\lr{TitForTat}         & $34.3$  & $2.7$   & $34.1$  & $9.7$   & $34.6$  & $\mathbf{37.3}$\\
\lr{Hawk}              & $-12.1$ & $-58.7$ & $12.0$  & $4.7$   & $12.1$  & $\mathbf{19.7}$\\
\lr{MasterAgent}       & $70.1$  & $-2.7$  & $96.8$  & $44.4$  & $96.0$  & $20.8$\\
\lr{Dove}              & $144.4$ & $135.9$ & $104.9$ & $64.6$  & $105.0$ & $\mathbf{113.1}$\\
\lr{BayesianThreshold} & $77.6$  & $9.4$   & $60.9$  & $15.0$  & $63.5$  & $24.7$\\
\midrule
\textbf{امتیاز}        & \multicolumn{2}{c}{$52.9$}
                       & \multicolumn{2}{c}{$56.8$}
                       & \multicolumn{2}{c}{$\mathbf{58.8}$}\\
\bottomrule
\end{tabular}
\caption{هر سه ستون عددهای \emph{واقعی} برد هستند. صندلی آمریکا مقابل
\lr{Dove} از $64.6$ به $113.1$ رفت — همان ضعفی که تشخیص داده بودیم.}
\label{tab:board3}
\end{table}

\textbf{اول، تشخیص درست بود.} صندلی آمریکا مقابل \lr{Dove} از $64.6$ به
$113.1$ و مقابل \lr{TitForTat} از $9.7$ به $37.3$ رفت. حریف صبور دقیقاً
همان چیزی بود که در استخر کم داشتیم.

\textbf{دوم، بازسازی برای سومین بار روی چهار حریف دقیق ماند.} تخمین ما
برای آن چهار حریف $64.8$ و $35.0$ بود و برد $64.8$ و $35.0$ داد. یعنی
دو سومِ معیار را با اطمینان کامل در اختیار داریم.

\textbf{سوم، و مهم‌تر از هر دو: برآوردگرمان زیادی خوش‌بین بود و خودش
هشدارش را داده بود.} تخمین $64.8$ بود و واقعیت $58.8$. با سه لنگر، حالا
می‌شود دید چرا:

\begin{table}[h]\centering\small
\begin{tabular}{lrrr}
\toprule
 & آپلود دوم & آپلود اول & آپلود سوم\\
\midrule
امتیاز جانشین صبور (صندلی آمریکا) & $-18.1$ & $9.2$  & $37.2$\\
امتیاز واقعیِ دو حریف نامعلوم     & $39.8$  & $72.7$ & $68.9$\\
\bottomrule
\end{tabular}
\caption{رابطه اول تند بالا می‌رود و بعد \emph{صاف می‌شود}. برازش خطیِ
دو-لنگره $106.2$ پیش‌بینی می‌کرد، واقعیت $68.9$ بود — و همان اختلاف،
دقیقاً همان $6$ واحد خطای تخمین کل است.}
\end{table}

\begin{findbox}
\textbf{یافته‌ی کلیدی ۱۱ — برون‌یابی را با درون‌یابی اشتباه نگیرید.}
برآوردگر ما کنار عددش نوشته بود «این کاندیدا $27.9$ واحد بیرون از بازه‌ی
لنگرهاست، پس این برون‌یابی است نه درون‌یابی؛ آن را یک \emph{جهت} بدانید
نه یک \emph{عدد}». جهت درست بود ($+2.0$ واحد بهبود واقعی) و عدد غلط
($+8.0$ ادعا شده). حالا برآوردگر بین لنگرها درون‌یابی تکه‌خطی می‌کند و
بیرون از بازه \emph{از برون‌یابی امتناع می‌کند}.
\end{findbox}

پیامد عملی این سقف مهم است: امتیاز آن دو حریف نامعلوم در صندلی آمریکا
حدود $70$ اشباع می‌شود و آپلود سوم با $68.9$ عملاً به آن رسیده است. پس از
«صبورتر شدن» چیز زیادی باقی نمانده؛ باقی‌مانده‌ی راه در همان چهار حریف
دقیق است.

\paragraph{چقدر راه باقی است، و کجا.}
با سه خواندن، می‌شود برای هر حریف بهترین عددی را که تا حالا گرفته‌ایم
کنار هم گذاشت. این یک \emph{کران بالای} قابل‌دستیابی نیست، ولی نشان
می‌دهد امتیاز کجا خوابیده است:

\begin{table}[h]\centering\small
\begin{tabular}{lrrr}
\toprule
حریف & ایران & آمریکا & معادل امتیاز کل\\
\midrule
\lr{Dove}              & $+39.4$ & $+22.8$ & $+5.18$\\
\lr{Random}            & $+4.9$  & $+51.7$ & $+4.72$\\
\lr{MasterAgent}       & $+0.8$  & $+23.6$ & $+2.03$\\
\lr{BayesianThreshold} & $+14.1$ & $+0.0$  & $+1.17$\\
\lr{TitForTat}, \lr{Hawk} & $0.0$ & $0.0$ & $0.00$ \quad(اشباع)\\
\bottomrule
\end{tabular}
\caption{فاصله‌ی آپلود سوم تا بهترین عددی که در هر خانه تا حالا گرفته‌ایم.
مجموع $13.1$ واحد است.}
\end{table}

بزرگ‌ترین قلم، صندلی آمریکا مقابل \lr{Random} است: از $113.8$ به $62.1$
افتاده. این همان بیش‌تصحیحی است که انتظارش را داشتیم — استخر
\lr{board+} سه حریف از هشت را «هیچ‌وقت نمی‌رود» گذاشت، در حالی که خودِ
برد دو از شش دارد، و عامل نتیجه گرفت که خروج زودهنگام همیشه امن است.
نسخه‌ی متعادل‌شده (\lr{--pool board+p}: فقط \lr{Patient} اضافه، دو از
شش) دقیقاً برای برگرداندن همین است.

\subsection{آخرین آزمایش: تیپ حریف در برابر رفتار حریف}
\label{sec:behav}

جدول بالا یک تناقض ظاهری دارد که ارزش دنبال کردن داشت. صندلی آمریکای ما
مقابل \lr{Random} در آپلود دوم $113.3$ گرفت و در آپلود سوم $62.1$؛ مقابل
\lr{Dove} دقیقاً برعکس ($64.6$ و $113.1$). دو بار ترکیب استخر را عوض
کردیم و هر بار عامل فقط \emph{یکی} از این دو رفتار را برداشت.

وقتی دو تغییر مستقل در توزیع آموزش نتواند دو رفتار را هم‌زمان نگه دارد،
احتمالاً مشکل در توزیع نیست. به ویژگی‌ها نگاه کردیم: هر ۱۶ ویژگی خلاصه‌ی
پسین ما روی \emph{تیپ ساختاری} حریف‌اند، $\theta_j=(\bar\rho_j,\bar m_j)$
— یعنی چقدر تاب‌آوری دارد و با چه سرعتی می‌سوزاند. ولی \lr{Random} و
\lr{Patient} از هم \emph{تیپ} متفاوتی ندارند؛ \emph{قاعده}شان فرق دارد، و
تفاوتشان تقریباً تماماً در \emph{نوسان} سیگنال است — چیزی که هیچ‌کدام از
آن ۱۶ ویژگی اندازه نمی‌گیرد.

سه آماره‌ی خام از جریان سیگنال حریف اضافه کردیم
(\lr{my\_team/theory\_behav.py}، ۱۹ ویژگی): انحراف معیار سیگنال، بیشینه‌ی
سیگنال تا اینجا، و کسر مراحلی که ساکت مانده. عمداً آماره‌ی خام‌اند نه
کمیت پسین: نظریه‌ای اضافه نمی‌کنند، فقط همان محوری را اضافه می‌کنند که
نظریه پوشش نمی‌دهد. پیش از صرف یک ساعت آموزش، بررسی کردیم که واقعاً جدا
می‌کنند یا نه:

\begin{table}[h]\centering\small
\begin{tabular}{lrrr}
\toprule
حریف & انحراف معیار & بیشینه & کسر سکوت\\
\midrule
\lr{Random}            & $0.159$ & $0.642$ & $0.225$\\
\lr{Patient}, \lr{Dove} & $0.000$ & $0.000$ & $\mathbf{1.000}$\\
\lr{Hawk}              & $0.000$ & $1.000$ & $0.000$\\
\lr{TitForTat}         & $0.325$ & $0.750$ & $0.250$\\
\lr{Plateau}           & $0.108$ & $0.500$ & $0.000$\\
\bottomrule
\end{tabular}
\caption{میانگین سه ویژگی تازه پس از ۵ مرحله، روی ۳۰ اپیزود. دقیقاً همان
دو حریفی که عامل قاطی‌شان می‌کرد — \lr{Random} و \lr{Patient} — بیشترین
فاصله را دارند.}
\end{table}

چهار بذر آموزش دادیم. نتیجه \emph{نصفه} بود، و نصفه‌اش هم دقیقاً آن نصفی
نبود که هدف گرفته بودیم:

\begin{table}[h]\centering\small
\begin{tabular}{lrrr}
\toprule
خانه (روی حریف‌های دقیق) & آپلود دوم & آپلود سوم & ۱۹-ویژگی\\
\midrule
\lr{Random} — آمریکا     & $\mathbf{113.3}$ & $62.1$ & $64.1$\\
\lr{MasterAgent} — ایران & $\mathbf{96.8}$  & $96.0$ & $86.6$\\
\lr{TitForTat} — ایران   & $34.1$ & $34.6$ & $\mathbf{41.1}$\\
\lr{Hawk} — ایران        & $12.0$ & $12.1$ & $\mathbf{22.8}$\\
\midrule
امتیاز کل (واقعی / تخمین) & $56.8$ & $58.8$ & $\approx59.8$\\
\bottomrule
\end{tabular}
\end{table}

\begin{findbox}
\textbf{یافته‌ی کلیدی ۱۲ — دادنِ یک محور به شبکه، تضمین نمی‌کند از آن
استفاده کند.} تشخیص درست بود: ویژگی‌ها واقعاً \lr{Random} را از
\lr{Patient} جدا می‌کنند، و جدول بالا این را پیش از آموزش نشان داد. ولی
هدف اصلی — برگرداندن $113.3$ در خانه‌ی \lr{Random}/آمریکا — \emph{حاصل
نشد} ($64.1$). سود از جای دیگری آمد: صندلی ایران مقابل \lr{TitForTat} و
\lr{Hawk} به بهترین عدد کل پروژه رسید. یعنی بازنمایی یک \emph{شرط لازم}
بود، نه کافی؛ سیاستِ صندلی آمریکا به دلیلی مستقل از قابلیت تشخیص، حاضر
نیست تعهد بالا را نگه دارد.
\end{findbox}

به همین دلیل چک‌پوینت تحویلی همان نسخه‌ی ۱۶-ویژگی‌ای است: بهبود
$\approx1$ واحدیِ نسخه‌ی ۱۹-ویژگی‌ای \emph{تخمین} است، در حالی که $58.8$
یک عدد \emph{اندازه‌گیری‌شده} است. نسخه‌ی ۱۹-ویژگی‌ای در بسته می‌ماند
(\lr{my\_team/theory\_behav.py}) چون آزمایشش و نتیجه‌ی نصفه‌کاره‌اش خودش
بخشی از یافته‌هاست.

\subsection{جست‌وجو برای عبور از $61.7$، و اینکه چرا متوقفش کردیم}
\label{sec:plateau}

آپلود چهارم $61.7$ گرفت. بعد از آن دوازده چک‌پوینت دیگر آموزش دادیم و
\emph{هیچ‌کدام} بهتر نشد. چون این بخش عمده‌ی کار پایانی پروژه است، کامل
گزارشش می‌کنیم؛ نتیجه‌ی منفی هم نتیجه است.

\begin{table}[h]\centering\small
\begin{tabular}{llr}
\toprule
تغییری که آزمودیم & منطقش & بهترین تخمین\\
\midrule
تعادل استخر (\lr{board+p}) & \lr{board+} زیادی صبور بود & $59.2$\\
آموزش $1.4$ میلیون قدمی & شاید همگرا نشده & متوقف شد\\
ویژگی «چقدر شاهد داریم» & $sd=0$ برای «ثابت» و برای & $59.7$\\
 & «یک مشاهده» یکی بود & \\
مواجهه‌ی اسکریپتی $56\%\to85\%$ & سیاست شرطی داده‌ی بیشتری & $61.0$\\
 & لازم دارد & \\
تنظیم میانی ($p_{\text{self}}=0.22$) & نگه داشتن هر دو خاصیت & $59.6$\\
\textbf{بذرهای تازه از همان دستور} & \textbf{کاهش نوفه} & $59.4$\\
\bottomrule
\end{tabular}
\caption{هیچ‌کدام به $61.7$ نرسید.}
\end{table}

توزیع خودِ چهارده چک‌پوینتِ خانواده‌ی ۱۹-ویژگی‌ای توضیح می‌دهد چرا:

\begin{center}\small
کمینه $51.5$ \quad میانه $58.7$ \quad میانگین $58.6$ \quad
انحراف معیار $2.33$ \quad بیشینه $\mathbf{61.7}$
\end{center}

\begin{findbox}
\textbf{یافته‌ی کلیدی ۱۳ — بخش زیادی از آن $61.7$، قرعه‌ی بذر بود.}
چک‌پوینت ارسالی $1.5$ انحراف معیار بالای میانگین بقیه‌ی خانواده‌ی خودش
است، و هشت بذر دیگر از \emph{همان دستور دقیق} هیچ‌کدام به آن نرسیدند
(بهترینشان $59.4$). یعنی $61.7$ کف توانایی این روش نیست، سقفِ خوش‌شانسیِ
آن است. با انحراف معیار $2.33$، فاصله‌ی $61.7$ تا میانه‌ی $58.7$ حدود
$1.3\sigma$ است — چیزی که با تکرار به‌دست می‌آید، نه با فهم بهتر.
\end{findbox}

یک تصحیح صادقانه هم لازم است. وسط همین جست‌وجو، \lr{by\_s2} خانه‌ی
\lr{Random} در صندلی آمریکا را از $64.1$ به $80.4$ برد و نرخ خروج ما را
از $70\%$ به $37.5\%$ رساند، و ما این را «جهت اثبات‌شده» نامیدیم. چهار
بذر بعدی با همان تنظیم آن را \emph{بازتولید نکردند} (همه حدود $35.5$).
پس آن هم یک قرعه بود، نه خاصیت تنظیم. این پنجمین فرضیه‌ای است که در این
پروژه با داده رد شد، و ثبتش می‌کنیم چون الگویش خودش درس است: در یک
جست‌وجوی پرنوفه، \emph{یک} مشاهده‌ی خوب تقریباً هیچ‌وقت شاهد کافی نیست.

نتیجه‌ی عملی: کار را همین‌جا متوقف کردیم. برای عبور مطمئن از $61.7$ یا
باید ده‌ها بذر دیگر زد (که یافته‌ی علمی‌ای تولید نمی‌کند) یا روش
اساساً بهتری داشت که در زمان باقی‌مانده پیدایش نکردیم.

\subsection{اجرای رسمی همان اجرای تمرینی نیست}
\label{sec:official}

آخرین چیزی که باید صریح گفته شود این است که جدول تمرینی، جدول نمره نیست.
\lr{docs/DESIGN.md} بخش ۱۰ دو تفاوت را مشخص می‌کند:

\begin{itemize}[leftmargin=1.5em,itemsep=3pt]
  \item \textbf{بند ۱۰-۳:} در اجرای رسمی «هر شرکت‌کننده با هر
        شرکت‌کننده‌ی \emph{دیگر} و با هر پنج حریف اسکریپتی بازی می‌کند، در
        هر دو نقش، روی یک فهرست بذر \emph{مشترک}». یعنی استخر هم عوض
        می‌شود، نه فقط بذر: به‌جای «۵ اسکریپتی + بنچ‌مارک استاد»،
        می‌شود «۵ اسکریپتی + ایجنت همه‌ی هم‌کلاسی‌ها».
  \item \textbf{بند ۱۰-۴:} اجرای نهایی با \lr{--seed0} متفاوت و
        اعلام‌نشده انجام می‌شود.
\end{itemize}

نکته‌ی دلگرم‌کننده در بند ۱۰-۳ این است که فهرست بذر بین همه‌ی جفت‌ها
مشترک است، پس همه دقیقاً همان قرعه‌ها را می‌بینند و شانس بذر به کسی
امتیاز نمی‌دهد — خودِ سند می‌گوید «تفاوت مهارت، نه شانس بذر، رتبه را
می‌سازد». برای اینکه این را با عدد بسنجیم، امتیاز را روی همان چهار
حریفی که دقیقاً بازتولید می‌شوند، روی پنج پایه‌ی بذر مجزا حساب کردیم
(\lr{analysis/seed\_shift.py}):

\begin{table}[h]\centering\small
\begin{tabular}{lrrrrrrr}
\toprule
چک‌پوینت & $20260901$ & $31000000$ & $31104729$ & $31209458$ & $31314187$
         & میانگین & دامنه\\
\midrule
آپلود اول  & $33.4$ & $32.8$ & $30.9$ & $36.8$ & $33.2$ & $33.4$ & $5.9$\\
آپلود دوم  & $\mathbf{54.6}$ & $52.0$ & $52.2$ & $51.8$ & $51.2$
           & $\mathbf{52.3}$ & $\mathbf{3.3}$\\
\bottomrule
\end{tabular}
\caption{فقط روی چهار مرجع دقیق، تا خطای بازسازی وارد این سنجش نشود.
پایه‌ی $20260901$ همان بذر تمرینی منتشرشده است.}
\end{table}

دو چیز از این جدول درمی‌آید. اول اینکه بذر منتشرشده حدود $2.3$ واحد به
نفع ماست، پس روی بذر تازه انتظار عدد کمی پایین‌تر منطقی است. دوم اینکه
دامنه‌ی نوسان روی بذرهای دیده‌نشده فقط $3.3$ واحد است و چک‌پوینت دوم از
اولی \emph{پایدارتر} هم هست ($3.3$ در برابر $5.9$) — یعنی برتری‌اش خاصیت
سیاست است نه شانس فهرست بذر.

ریسک واقعی جای دیگری است: حریف‌های تازه ایجنت‌های \emph{یادگرفته‌ی}
هم‌کلاسی‌ها هستند. چک‌پوینت فعلی دقیقاً مقابل همین تیپ قوی است
(\lr{Hawk} از $-58.7$ به $+4.7$ و \lr{MasterAgent} از $-2.7$ به $44.4$)،
ولی ضعفش مقابل حریف صبورِ منفعل هم همان‌جا ظاهر می‌شود اگر کسی چنین
ایجنتی بفرستد. برای همین تصحیح استخر، فارغ از عددِ جدول تمرینی، کار
درستی است.

% ==================================================================== %
\section{تفسیر نتایج}\label{sec:interp}
% ==================================================================== %

\begin{itemize}[leftmargin=1.5em,itemsep=4pt]
  \item \textbf{چرا تقابل پرهزینه ادامه پیدا می‌کند؟}
        چون آستانه‌ی ادامه پایین است. با $B\approx150$ و $X\approx26$ و
        هزینه‌ی جاری $\approx1.1$، بازیکن باید ادامه دهد مگر آنکه شانسش زیر
        حدود $1\%$ باشد. دو آستانه‌ی این‌قدر پایین به‌راحتی هم‌زمان ارضا
        می‌شوند — پس \emph{ادامه‌ی دوطرفه یک تعادل است، نه یک خطا}.
  \item \textbf{چرا افزایش هزینه لزوماً تقابل را تمام نمی‌کند؟}
        چون هزینه از کانال $R/m$ وارد می‌شود و در گزاره~\ref{pr:gate} فقط در
        \emph{صورت} کسر است. افزایش $R$ تا وقتی $q$ را زیر $q^\ast$ نبرد
        رفتار را عوض نمی‌کند. علاوه بر این فشار بیشتر $\tilde\sig_j$ را بالا
        می‌برد که $X_i$ را هم پایین می‌آورد — یعنی هم‌زمان که ماندن را
        گران‌تر می‌کند، رفتن را هم بی‌آبروتر می‌کند.
  \item \textbf{چرا مدیریت هزینه از خود هزینه مهم‌تر است؟}
        چون $m$ در \emph{مخرج} است و بازخورد دارد
        ($m=\mbar(0.35+0.65\ehat)$). طبق گزاره~\ref{pr:collapse}،
        $T_{\text{فروپاشی}}$ نسبت به $\mbar$ خطی و نسبت به $R$ وارون است، و
        نسبت به تاب‌آوری باقی‌مانده محدب — یعنی ضعف مدیریتی خودش را تشدید
        می‌کند.
  \item \textbf{چرا باورها از واقعیت مادی مهم‌تر می‌شوند؟}
        چون در گزاره~\ref{pr:gate} تنها چیزی که ضرب در $B-X$ می‌شود $q$ است،
        و $q$ یک \emph{باور} است. دو بازیکن با وضعیت مادی یکسان ولی باور
        متفاوت، رفتار کاملاً متفاوت دارند.
  \item \textbf{چرا سیگنال قوی همیشه بهترین نیست؟}
        چون هزینه‌اش درجه دو است، $K$ را بالا می‌برد و $X$ را می‌خورد، و
        پنجره‌ی میانجی‌گری را می‌بندد. سنجش ما نشان داد ترتیب بهینه
        $0.75>1.0>0.0>0.5$ است.
  \item \textbf{چرا خروج باید از نظر سیاسی امکان‌پذیر شود؟}
        چون $+20M$ بزرگ‌ترین جمله‌ی $X$ است و در داده، خروج‌ها با پنجره‌ی
        باز بیشتر رخ می‌دهند. \emph{ساختن} مسیر خروج تقریباً تنها راه
        پایان کم‌ضرر است — و برعکسش هم در داده دیده می‌شود: عاملی که
        دائم تشدید می‌کند خودش پنجره را بسته نگه می‌دارد و آن راه را روی
        خودش می‌بندد.
  \item \textbf{چرا حفظ انعطاف‌پذیری بخشی از عقلانیت راهبردی است؟}
        چون $K$ ماندگار است ($0.75$ فراموشی در هر مرحله) و $X$ را خطی
        می‌خورد. عامل آموزش‌دیده هم همین را نشان می‌دهد: با $K$ بالا احتمال
        خروج داوطلبانه‌اش عملاً به صفر می‌رسد — یعنی واقعاً \emph{در تعهدات
        خودش گیر می‌افتد}.
\end{itemize}

% ==================================================================== %
\section{محدودیت‌های مدل}\label{sec:limits}
% ==================================================================== %

\paragraph{چه چیزهایی ساده‌سازی شده؟}
تقابل به دو بازیگر یکپارچه، یک بعد سیگنال، و یک متغیر حالت اسکالر
(تاب‌آوری) ساده شده.

\paragraph{آیا دو بازیکنی دیدن کافی است؟}
نه. بازیگران منطقه‌ای، متحدان و نهادهای بین‌المللی فقط از راه یک متغیر
دوحالتی $M^t$ وارد شده‌اند که \emph{از بیرون} تعیین می‌شود و به هیچ منافع مستقلی پاسخ
نمی‌دهد. در واقعیت میانجی خودش بازیگر است.

\paragraph{افکار عمومی و نهادهای داخلی؟}
کل سیاست داخلی در $\alpha_i K_i$ خلاصه شده — یک تابع خطی از تعهدات عمومی
خود بازیگر. هیچ تفکیکی بین جناح‌ها، هیچ چرخه‌ی انتخاباتی، و هیچ امکان
تغییر حکومت نیست.

\paragraph{آیا فرض عقلانیت کامل قابل دفاع است؟}
تا حدی. ما \emph{عقلانیت ترتیبی} را فرض کردیم ولی آنچه واقعاً به دست آمد یک
تعادل \emph{تقریبی} حاصل از یادگیری است. جالب اینکه همین امر واقع‌گرایانه‌تر
هم هست: یک عامل با بودجه‌ی محاسباتی محدود که سیاستش را از تجربه می‌سازد،
شبیه‌تر به تصمیم‌گیر واقعی است تا یک حل‌کننده‌ی کامل \lr{PBE}.

\paragraph{آیا باورها همیشه با قاعده‌ی بیز به‌روز می‌شوند؟}
در نظریه بله، در پیاده‌سازی \emph{تقریباً}. فیلتر ما نویز‌ی $\xi$ را در
برآورد نقطه‌ای نادیده می‌گیرد و کانال سیگنال از یک درست‌نمایی \emph{پاسخ کوانتال}
استفاده می‌کند نه استراتژی تعادلی واقعی حریف. جدول اعتبارسنجی نشان می‌دهد
این تقریب‌ها بایاس معناداری وارد نمی‌کنند ($+0.0033$) ولی خطای پراکندگی
باقی می‌ماند.

\paragraph{آیا سیگنال‌ها همیشه قابل مشاهده و تفسیرند؟}
در مدل بله — تاریخچه‌ی عمومی \emph{کامل} است. این یک ساده‌سازی جدی است:
در واقعیت اقدامات مبهم‌اند، دیر مخابره می‌شوند و می‌توانند اشتباه تفسیر شوند.
مدل \emph{عدم‌قطعیت درباره‌ی نوع} دارد ولی \emph{عدم‌قطعیت درباره‌ی مشاهده}
ندارد.

\paragraph{آیا مدل قدرت پیش‌بینی دارد؟}
به نظر ما نه، و ادعایش را هم نداریم. مدل یک \emph{چارچوب تحلیلی} است:
می‌گوید کدام کمیت‌ها با هم رابطه دارند و در کدام جهت. حتی همین حرف هم مشروط
است، چون بخش زیادی از نتایج عددی به مقادیر تنظیم‌شده‌ی
\lr{see/config.py} حساس است.

\paragraph{کدام نتیجه به تغییر مفروضات حساس‌تر است؟}
به‌ترتیب حساسیت:
\begin{enumerate}[leftmargin=1.7em,itemsep=2pt]
  \item \textbf{«خروج داوطلبانه بهینه نیست»} کاملاً به نسبت
        $b_0/x_0=175/26\approx6.7$ وابسته است. سند \lr{DESIGN.md} خودش
        می‌گوید $b_0$ از $110$ به $175$ و $x_0$ از $46$ به $26$ تغییر کرده
        تا بازی بی‌معنی نشود. با مقادیر اولیه، نتیجه \emph{برعکس} می‌شد.
  \item \textbf{«$\sig=0.75$ بهترین نقطه است»} به دو جزئیات پیاده‌سازی
        وابسته است: آستانه‌ی \lr{HIGH\_SIGNAL} برابر $0.75$ و معافیت $\varphi$
        در \lr{HOLD}. این دومی در متن اصلی اصلاً نیست و — همان‌طور که در
        بخش~\ref{sec:credflex} گفتیم — تا حدی خلاف قصد متن عمل می‌کند.
  \item \textbf{نتایج تفکیک نوع} به شکاف $7.2$ برابری $\varphi$ وابسته‌اند
        که خودش از $\varphi_0=1.2$ و کف‌های $0.4$ و $0.5$ می‌آید.
  \item کم‌حساس‌ترین نتیجه \emph{ساختار آستانه‌ای} (گزاره~\ref{pr:cutoff})
        است، چون فقط به \emph{علامت} مشتق‌ها بستگی دارد نه به مقادیر — و
        همان علامت‌ها در متن اصلی به‌عنوان فرض آمده‌اند.
\end{enumerate}

\begin{keybox}
جمع‌بندی این بخش: مدل درباره‌ی \emph{جهت} روابط قابل دفاع است و درباره‌ی
\emph{اندازه}ی آن‌ها نه. هر عددی که در این گزارش آمده، عدد \emph{این}
نمونه‌سازی است، نه عدد جهان.
\end{keybox}

% ==================================================================== %
\section{جمع‌بندی}\label{sec:concl}
% ==================================================================== %

پرسش شروع این بود که چرا تقابل پرهزینه ادامه پیدا می‌کند. پاسخ این گزارش،
در یک خط: \textbf{چون آستانه‌ی ادامه پایین است و ادامه دادن، خودش آن آستانه
را پایین‌تر می‌آورد.}

سه حلقه‌ی خودتقویت‌کننده این را می‌سازند:
\begin{enumerate}[leftmargin=1.7em,itemsep=3pt]
  \item \textbf{حلقه‌ی تعهد:} تشدید $\Rightarrow$ $K$ بالا $\Rightarrow$
        $X$ پایین $\Rightarrow$ مخرج $B-X$ بزرگ‌تر $\Rightarrow$
        $q^\ast$ کوچک‌تر $\Rightarrow$ ادامه عقلانی‌تر.
  \item \textbf{حلقه‌ی شکنندگی:} فرسایش $\Rightarrow$ $m$ پایین
        $\Rightarrow$ $g$ بالا $\Rightarrow$ فرسایش سریع‌تر.
  \item \textbf{حلقه‌ی میانجی:} تشدید $\Rightarrow$ پنجره بسته می‌ماند
        $\Rightarrow$ $X$ پایین می‌ماند $\Rightarrow$ خروج غیرممکن‌تر.
\end{enumerate}

و تنها چیزی که هر سه حلقه را می‌شکند \emph{خویشتن‌داری متقابل} است، چون
هم‌زمان $K$ را می‌خواباند، سوخت را کم می‌کند و پنجره‌ی میانجی‌گری را باز
می‌کند. این دقیقاً همان نتیجه‌ای است که متن اصلی به زبان کیفی می‌گوید، و ما
دیدیم که هم از معادلات درمی‌آید و هم عامل آموزش‌دیده مستقلاً به آن می‌رسد.

بخش پیاده‌سازی درس جداگانه‌ی خودش را داشت، و آن درس درباره‌ی
\emph{اندازه‌گیری} است نه درباره‌ی بازی. ما معیار نمره‌دهی را باز کردیم و به
رابطه‌ی $J_I+J_U>R_1+R_2$ رسیدیم، برایش نظریه ساختیم، سه بار شکست خوردیم و
بار چهارم بردیم — و بعد معلوم شد کل آن کار روی استخر \emph{اشتباهی} انجام
شده بود. معیار واقعی شش حریف دارد نه دو تا، و مهم‌تر اینکه اصلاً نسبی
نیست: امتیاز \lr{TitForTat} عددی ثابت است که ما نمی‌توانیم تکانش دهیم، پس
تمام آن ساختمانِ «محروم‌کردن حریف» روی برد بی‌معنی بود
(بخش~\ref{sec:realboard}).

این پنج شکست را عمداً در گزارش نگه داشته‌ایم، چون هر کدام یک درس مستقل
دارند: جریمه‌ای که مهلت بدهد فقط فرار را یاد می‌دهد؛ اندازه‌گیری روی سیاست
ثابت، پیش‌بینی سیاست آموزش‌دیده نیست؛ آزمونی که در حالت پیش‌پاافتاده همیشه
درست جواب بدهد اصلاً آزمون نیست؛ آماره‌ای که خروج اجباری را انتخاب حساب
کند علامت نتیجه را برعکس می‌کند؛ و — گران‌ترینشان — بهینه‌سازی روی معیاری
که \emph{در دسترس} است به‌جای معیاری که \emph{نمره می‌دهد}.

% ==================================================================== %
\section{استفاده از هوش مصنوعی}
% ==================================================================== %

دستورالعمل این پروژه (بند ۱۱) استفاده از ابزارهای هوش مصنوعی را برای فهم
اولیه‌ی متن، ویرایش نگارشی، تولید کد اولیه و مرتب‌سازی گزارش مجاز دانسته و
فقط خواسته است که در صورت استفاده‌ی قابل‌توجه، دانشجو توضیح دهد از این ابزارها
برای چه کاری استفاده کرده است. کاری که من کردم دقیقاً داخل همین چهار مورد
است و در ادامه مورد به مورد توضیحش می‌دهم. \emph{هیچ‌کدام} از تحلیل‌های
بازی‌نظری بخش‌های ۳ تا ۱۰ (دروازه‌ی ادامه، ساختار آستانه‌ای، تحلیل باور،
تک‌تقاطعی، تعارض اعتبار و انعطاف، و گزاره‌های تحلیلی) پاسخ آماده‌ای نیست که
گرفته باشم؛ فرمول‌بندی و اثبات‌ها را خودم نوشته‌ام.

روند کار من این‌طور بود که مسئله را برای هوش مصنوعی شرح دادم و با کمک آن روی
مبحث مسلط شدم. عملاً قدم‌به‌قدم یاد گرفتم که برای هر بخش پروژه باید چه کار کرد
و بعد خودم انجامش دادم.

مهم‌ترین نکته درباره‌ی شکل کار این است که یک‌بار پرسیدن و یک‌بار جواب گرفتن
نبود؛ یک \emph{رفت‌وبرگشت} طولانی بود. من مدام ایده می‌دادم — «این عدد
مشکوک است، بیا اندازه‌اش بگیریم»، «شاید مشکل از استخر حریف باشد»، «این را
جدا برای هر صندلی حساب کن» — و برای هر ایده کد \emph{تازه} می‌زدیم، اجرا
می‌کردیم، به عددها نگاه می‌کردم، و بعد یا ایده را دور می‌ریختم یا اصلاحش
می‌کردم و دوباره می‌رفتیم. بیشتر فایل‌های پوشه‌ی \lr{analysis/} دقیقاً
همین‌طور به وجود آمدند: هرکدام‌شان جواب یک سوالی است که وسط کار برایم پیش
آمد و برای جواب دادن به آن سوال کد جدید نوشتیم. مثلاً
\lr{frontier.py} و \lr{classpool.py} و \lr{denial\_probe.py} هیچ‌کدام از
اول در برنامه نبودند؛ وقتی دیدم از \lr{TitForTat} عقبیم پرسیدم «دقیقاً کجا
عقبیم؟» و برای جواب دادن مجبور شدیم سه بار کد تازه بزنیم تا به تجزیه‌ی
بخش~\ref{sec:tourn} رسیدیم.

جزئیات مورد به مورد:

\begin{itemize}[leftmargin=1.4em,itemsep=3pt]
  \item \textbf{فهم متن:} مدل $\Gamma$ و رابطه‌ی اجزایش را اول با کمک هوش
        مصنوعی برای خودم روشن کردم، مخصوصاً تفاوت $\theta_i$ (نوع ساختاری
        ثابت) و $e_i^t$ (حالت جاری) و اینکه چرا شرط تک‌تقاطعی همان چیزی است
        که سیگنال را اطلاعاتی می‌کند.
  \item \textbf{کدها:} کد را با کمک هوش مصنوعی زدم ولی خودم اصلاح کردم.
        نسخه‌ی اول فیلتر باور خیلی کند بود چون برای هر ذره مسیر را
        شبیه‌سازی می‌کرد؛ خودم آن را با حل بستن انتگرال
        (گزاره~\ref{pr:collapse}) جایگزین کردم که هم سریع‌تر است و هم
        تحلیلی‌تر. همچنین بعد از دیدن خروجی \lr{probe.py} فهمیدم مشکل
        اصلی استخر حریف است و الگوها را خودم طراحی کردم.
  \item \textbf{تئوری‌ها:} کل تئوری‌ها را خودم بازبینی کردم. اثبات گزاره‌ی
        ساختار آستانه‌ای و مشتق‌گیری زمان فروپاشی را قدم‌به‌قدم چک کردم، و
        عددهای جدول‌ها ($7.2$ برابر شکاف $\varphi$، بازه‌های $B$ و $X$) را
        مستقیم از \lr{config.py} دوباره حساب کردم تا مطمئن شوم با متن
        می‌خوانند.
  \item \textbf{یافته‌های خلاف انتظار:} دو موردی که بیشترین وقت را گرفتند
        (تله‌ی ابرپارامتر \lr{train.py} و معافیت $\varphi$ در \lr{HOLD})
        هیچ‌کدام از خواندن متن پیدا نشدند، بلکه از مقایسه‌ی رفتار
        اندازه‌گیری‌شده با پیش‌بینی نظری بیرون آمدند. جست‌وجوی استراتژی‌ها را هم
        دقیقاً برای همین نوشتم که ببینم عدد واقعی کجاست.
  \item \textbf{ایده دادن و کد تازه زدن، پشت سر هم:} این بیشترین وقت کار
        بود. چند نمونه‌ی مشخص از این چرخه:
        \begin{itemize}[leftmargin=1.2em,itemsep=1pt]
          \item «چرا آمریکا همیشه تسلیم می‌شود؟» $\to$ \lr{probe.py} را
                نوشتیم $\to$ دیدم ۱۰۰٪ است $\to$ حدس زدم تله‌ی ابرپارامتر
                است $\to$ درست بود.
          \item «شاید آن هم بهترین پاسخ باشد» $\to$ حساب کردیم $\to$
                واقعاً بود، و تشخیص اولم غلط بود؛ در گزارش تصحیحش کردم.
          \item «کدام صندلی عقب است؟» $\to$ \lr{frontier.py} را ساختیم
                $\to$ معلوم شد معیار به دو مسئله‌ی مستقل تجزیه می‌شود.
          \item «آیا این فقط اثر استخر کوچک است؟» $\to$ \lr{classpool.py}
                $\to$ نه، در همه‌ی اندازه‌ها هست.
          \item «اگر تسلیم نشویم چه؟» $\to$ \lr{denial\_probe.py} $\to$
                عدد نشان داد می‌ارزد، و تازه بعد از دیدن آن عدد بود که
                کانال \lr{(S4)} را به نظریه اضافه کردیم.
          \item بعد از اولین بارگذاری روی برد: «چرا عددِ برد با عدد ما
                نمی‌خواند؟» $\to$ متن خودِ برد را خواندم و دیدم شش حریف
                دارد نه دو تا $\to$ گفتم «پس بیا کل برد را محلی بسازیم»
                $\to$ \lr{board.py} و \lr{reference\_pool.py} را از صفر
                نوشتیم $\to$ اولین اجرا خیلی پرت بود، من حدس زدم شاید هر
                عامل با خودش هم بازی می‌کند $\to$ امتحان کردیم و خطا
                نصف شد.
          \item «چطور مطمئن شویم بازسازی درست است؟» $\to$ فکر کردم سه
                حریف را دقیق داریم، پس اگر آن سه‌تا دقیق دربیایند
                هارنس درست است $\to$ \lr{board\_detail.py} را برای همین
                نوشتیم $\to$ هر سه دقیق درآمدند و \lr{Hawk} هم به‌طور
                غیرمنتظره دقیق شد.
          \item «چرا مقابل \lr{Hawk} این‌قدر بد بازی می‌کنیم؟» $\to$
                نگاه کردم دیدم جریمه‌ی تسلیم دارد ما را وادار به جنگیدن
                در مسابقه‌ای می‌کند که باختش قطعی است $\to$ اول ضریب
                $\hat q$ را اضافه کردیم، و بعد فهمیدم اصلاً کل آن جمله
                برای معیار برد بی‌معنی است چون امتیاز حریف ثابت است.
          \item بعد از آپلود دوم: «تخمین ما $67$ بود و برد $56.8$ داد، پس
                کجا اشتباه است؟» $\to$ گفتم برازش دوباره‌ی \lr{Dove} را
                امتحان کنیم چون حالا دو آپلود داریم $\to$
                \lr{fit\_refs2.py} $\to$ \emph{هیچ} پارامتری هر دو را
                نساخت. اینجا بود که گفتم «پس بیا به‌جای حریف، رفتار خودمان
                را اندازه بگیریم» $\to$ \lr{sigprofile.py} $\to$ معلوم شد
                عامل ۹۶٪ مراحل $\sigma=0.75$ را نگه می‌دارد و هیچ‌وقت
                خارج نمی‌شود.
          \item «اگر نمی‌توانیم بازسازی‌اش کنیم، شاید حریفی داشته باشیم که
                \emph{ترتیبش} را درست بچیند» $\to$ \lr{proxy\_probe.py}
                $\to$ \lr{Patient} تنها موردی بود که این کار را می‌کرد، و
                همان لحظه فهمیدم چرا از استخر افتاده بود.
          \item «چرا خانه‌ی \lr{Random} افت کرد؟» $\to$ \lr{seat\_probe.py}
                را نوشتیم تا رفتار را به تفکیک (چک‌پوینت، صندلی، حریف)
                ببینیم $\to$ دیدم عامل ۷۰\% مواقع در مرحله‌ی $\sim2$ خودش
                خارج می‌شود در حالی که \lr{Random} داشت اول می‌رفت. حدس
                زدم چون شاهد کم است \lr{Random} شبیه \lr{Hawk} دیده
                می‌شود $\to$ ویژگی شمارش شاهد را اضافه کردیم $\to$
                \emph{جواب نداد}، و در گزارش به‌عنوان فرضیه‌ی رد شده
                ثبتش کردم.
          \item «آیا $61.7$ مهارت است یا شانس؟» $\to$ هشت بذر دیگر از
                همان دستور زدیم $\to$ هیچ‌کدام نرسید، و توزیع نشان داد
                $1.5\sigma$ بالای میانگین است (یافته‌ی کلیدی ۱۳). این
                همان جایی بود که تصمیم گرفتم جست‌وجو را متوقف کنم.
          \item وقتی استاد پوشه‌ی \lr{tournament/} را منتشر کرد: «مطمئن
                شویم چیزی که ما داشتیم همان است» $\to$
                \lr{verify\_tournament.py} $\to$ اول \lr{diff} گفت کاملاً
                متفاوت‌اند و یک لحظه نگران شدم؛ بعد دیدم فقط
                \lr{CRLF}/\lr{LF} است و پس از یکسان‌سازی بایت‌به‌بایت
                یکی‌اند.
        \end{itemize}
  \item \textbf{نگارش:} متن‌ها را خودم نوشتم ولی با کمک هوش مصنوعی
        ویرایششان کردم و در نهایت تبدیل به لاتک برای گزارش کردم.
\end{itemize}

مسئولیت درستی مفروضات، نمادها، روابط، گزاره‌های تحلیلی و تفسیر نتایج با
خودم است و می‌توانم در جلسه‌ی دفاع از آن‌ها دفاع کنم.

% ==================================================================== %
\appendix
\section{راهنمای کد و بازتولید}
% ==================================================================== %

\begin{table}[h]\centering\small
\begin{tabular}{ll}
\toprule
فایل & کار\\
\midrule
\lr{my\_team/theory.py}          & نظریه‌ی ارسالی: فیلتر بیزی + شکل‌دهی پتانسیل‌پایه\\
\lr{my\_team/opponents.py}       & پنج الگوی حریف برای آموزش\\
\lr{my\_team/train\_plus.py}     & مربی با استخر حریف غنی‌شده\\
\lr{see/tournament/runner.py}    & ماژول اصلی تورنمنت (از تاریخچه‌ی گیت)\\
\lr{analysis/filter\_check.py}   & اعتبارسنجی فیلتر در برابر مقدار واقعی\\
\lr{analysis/probe.py}           & کاوش رفتاری یک چک‌پوینت\\
\lr{analysis/sweep.py}           & جست‌وجوی فضای استراتژی\\
\lr{analysis/propositions.py}    & آزمون چهار گزاره‌ی تحلیلی (بخش ۵-۵)\\
\lr{analysis/compare.py}         & تورنمنت دوره‌ای همه‌ی پیکربندی‌ها\\
\lr{analysis/frontier.py}        & تجزیه‌ی معیار به دو صندلی، رابطه‌ی~\eqref{eq:crit}\\
\lr{analysis/classpool.py}       & همان معیار در استخرهای بزرگ‌تر\\
\lr{analysis/diagnose\_iran.py}  & تفکیک شکاف صندلی ایران حریف‌به‌حریف\\
\lr{analysis/denial\_probe.py}   & سنجش سود و هزینه‌ی کانال \lr{(S4)}\\
\lr{analysis/seat\_split.py}     & تفکیک $J_I$ و $J_U$ برای هر چک‌پوینت\\
\lr{analysis/blanket\_cost.py}   & هزینه‌ی «نرفتن» حریف‌به‌حریف\\
\lr{analysis/seed\_robust.py}    & برتری روی شش بلوک بذر مجزا\\
\lr{analysis/splice.py}          & ساخت چک‌پوینت از بهترین صندلی هر طرف\\
\lr{analysis/k\_control.py}      & گزاره‌ی ۴ با کنترل مرحله و تاب‌آوری\\
\lr{analysis/benchmark\_table.py}& جدول ارزیابی همراه بنچمارک استاد\\
\lr{my\_team/reference\_pool.py}  & بازسازی شش حریف مرجع برد کلاسی\\
\lr{analysis/board.py}           & بازتولید محلی برد و برازش حریفان\\
\lr{analysis/board\_detail.py}    & امتیاز یک چک‌پوینت حریف‌به‌حریف روی برد\\
\lr{analysis/fit\_refs.py}        & برازش \lr{Dove} و \lr{BayesianThreshold}\\
\lr{analysis/select\_final.py}   & انتخاب چک‌پوینت ارسالی روی معیار نمره‌دهی\\
\lr{my\_team/theory\_behav.py}    & نظریه‌ی ارسالی: ۱۹ ویژگی (\lr{+} رفتار حریف)\\
\lr{my\_team/theory\_behav2.py}   & نسخه‌ی ۲۰ ویژگی؛ آزمایشِ رد شده (بخش~\ref{sec:plateau})\\
\lr{analysis/fit\_refs2.py}       & برازش دوباره روی \emph{دو} آپلود؛ نشان داد شکست می‌خورد\\
\lr{analysis/sigprofile.py}       & رفتار یک چک‌پوینت مقابل حریف کاملاً منفعل\\
\lr{analysis/proxy\_probe.py}     & یافتن حریف جانشین که ترتیب برد را حفظ کند\\
\lr{analysis/seat\_probe.py}      & رفتار به تفکیک (چک‌پوینت، صندلی، حریف)\\
\lr{analysis/select2.py}          & انتخاب: خانه‌های دقیق جدا از جانشین\\
\lr{analysis/seed\_shift.py}      & حساسیت امتیاز به پایه‌ی بذر\\
\lr{analysis/verify\_tournament.py}& تطبیق \lr{see/tournament/} با نسخه‌ی منتشرشده\\
\bottomrule
\end{tabular}
\end{table}

\noindent\textbf{بازتولید نتیجه‌ی نهایی.} دستورهای زیر از ریشه‌ی
\lr{see-project} اجرا می‌شوند:

\begin{latin}\small
\begin{verbatim}
# train against the reconstructed board pool, self-play off
python my_team/train_plus.py --theory my_team/theory.py \
    --team "402111237" --steps 700000 --pool board \
    --p-self 0.0 --p-scripted 1.0 --entropy 0.02 \
    --entropy-final 0.004 --seed 94 --out submission.pt

# select on the board replica, not on evaluate.py
python analysis/board_detail.py submission.pt --episodes 40 --compare
python analysis/board.py --episodes 40 --self-play
\end{verbatim}
\end{latin}

\end{document}
